In this section, we state and prove the circle decomposition theorem, Theorem~\reffour{theorem:intro-F} from the introduction, which describes the behaviour of Verdier-localising functors under the operation $(\C,\QF) \mapsto (\C,\QF)_{\B\ZZ}$ of tensoring with the circle. Our framework reproduces both the Bass-Heller-Swan and the Ranicki-Shaneson decompositions in a unified manner. 
Our main result is stated and put into perspective in \S\reffour{subsection:shaneson} and proven in \S\reffour{subsection:proof}. In \S\reffour{subsection:splittings}, we revisit Ranicki's observation about the difference between the algebraic and geometric Ranicki-Shaneson decompositions in the setting of Poincaré $\infty$-categories.



\subsection{Statement of the circle decomposition theorem}
\label{subsection:shaneson}%


 


To explain the precise statement of the circle decomposition theorem, we recall that for a monoid $M$, we may form the pointed category $\B M$ with a single object whose endomorphisms are $M$, naturally in the monoid $M$. When $M$ is a group we will also view $\B M$ as a pointed space. The examples we have in mind are the monoid $\NN$ of natural numbers and the group $\ZZ$ of the integers.
In particular, 
for a stable $\infty$-category we may form $\C_{\B\NN}$ and $\C_{\B\ZZ}$, the tensor of $\C$ with $\B\NN$ and $\B\ZZ$, respectively. The monoid inclusion $\NN \hrar \ZZ$ then determines a map $q\colon\B\NN \to \B\ZZ$ and hence an exact functor $\C_{\B\NN} \to \C_{\B\ZZ}$ for every stable $\infty$-category $\C$. 
Similarly, we may form these tensors with $(\C,\QF)$ a hermitian $\infty$-category and if $(\C,\QF)$ is Poincar\'e, then so is $(\C,\QF)_{\B\ZZ}$, while $(\C,\QF)_{\B\NN}$ typically is not. Finally, we note that the functors $\ast \to \B\NN \to \B\ZZ \to \ast$ make $\C$ and $(\C,\QF)$ a retract of $\C_{\B\NN}$ and $(\C,\QF)_{\B\ZZ}$, respectively. To keep in mind the effect of tensoring with $\B\ZZ$, we recall the following.

\begin{example}
\label{example:tensor-BZ}%
Suppose $A$ is an $\Eone$-ring, $(M,N,\alpha \colon N \to M^{\tC})$ a module with genuine involution over $A$, and $(\Mod^\free(A),\QF^\alpha_M)$ the associated Poincar\'e $\infty$-category. Then we have 
\[
(\Mod^\free(A),\QF^\alpha_M)_{\B\ZZ} \simeq (\Mod^\free(A[\ZZ]),\QF^{\alpha[\ZZ]}_{M[\ZZ]})
\]
where $\alpha[\ZZ]$ is the map $N \simeq N \otimes \SS \to M^{\tC} \otimes \SS[\ZZ]^{\tC} \to M[\ZZ]^{\tC}$, see Example \refone{example:tensor-of-rings}. In particular, we have:
\begin{enumerate}
\item
\label{item:ex-tensor-BZ-genuine-family}%
If $(R,\sigma)$ is a ring with involution, then $(\Dfree(R),\QF^{\geq m}_\sigma)_{\B\ZZ} \simeq (\Dfree(R[\ZZ]),\QF^{\geq m}_{\bar{\sigma}})$, where $\bar{\sigma}$ is the unique extension of the involution $\sigma$ on $R$ which is the inversion map on $\ZZ$; under the canonical isomorphism $R[\ZZ] \cong R[t^{\pm 1}]$, this means $\bar{\sigma}(t) = t^{-1}$.
\item 
For a ring $R$ with form parameter $\lambda=[M_{\Ct} \to Q \to M^{\Ct}]$ and associated Poincar\'e structure $\QF^{\g\lambda}_M$, we have an equivalence
\[
(\Dfree(R),\QF_M^{\g\lambda})_{\mathrm{B}\mathbb Z}  \simeq (\Dfree(R[\ZZ]),\QF_{M[\ZZ]}^{\g\lambda[\ZZ]})
\]
where $\lambda[\ZZ]$ is the form parameter on $R[\ZZ]$ whose associated module with involution is $M[\ZZ] = \oplus_{n \in \ZZ}M$ as above and quadratic module given by $Q \oplus \big[\oplus_{n > 0}M\big]$. While the operation $\lambda \mapsto \lambda[\ZZ]$ preserves quadratic, even, and symmetric form parameters %
by \reffouritem{item:ex-tensor-BZ-genuine-family} above, there are other classes of form parameters which do not obey such a compatibility. For instance, if $(R,\sig)$ is a commutative ring with involution and we consider the Frobenius form parameter $\lambda = R^\fr$ as in Example~\reffour{example:form_parameter_over_Tate_Poincare}, then $\lambda[\ZZ]$ is not the Frobenius form parameter on $R[\ZZ]$, unless $R$ is zero. Otherwise put, on the level of associated Poincaré structures we have that 
$\QF^{\g\fr[\ZZ]}_{\sig[\ZZ]} \not{\simeq} \QF^{\g\fr}_{\sigma[\ZZ]}$ 
as Poincaré structures on $\Dfree(R[\ZZ])$ %
(indeed, the $\pi_0$ of the $R[\ZZ]$-complex representing the linear part is $R$ in the former case and $R[\ZZ]$ in the latter, see Example~\reffour{example:modules-over-genuine} below).
\item For a spherical fibration $\xi \colon X \rightarrow \Pic(\SS)$ on a space $X$ and the associated visible Poincar\'e category $(\Sp_X^\free,\QF^\vis_\xi)$ we have 
$(\Sp_X^\free,\QF^\vis_\xi)_{\B\ZZ} \simeq (\Sp_{X\times S^1}^\free,\QF^\vis_{p^*(\xi)})$
where $p\colon X \times S^1 \to X$ is the projection.
\end{enumerate}
\end{example}



\begin{construction}
\label{construction:nil}%
\label{remark:splitting}%
Let $\G\colon \Catx \to \A$ be a functor from $\Catx$ to a stable $\infty$-category $\A$. We define
$\N\G\colon \Catx \to \A$ to be the functor given by the cofibre of the map
$\G(\C) \to \G(\C_{\B \NN})$ or equivalently the fibre of the map $\G(\C_{\B\NN}) \to \G(\C)$ and obtain a natural splitting 
\[
\G(\C_{\B\NN}) \simeq \G(\C) \oplus \N\G(\C)
\]
and refer to it as the \defi{nil term} of $\G$. The terminology comes from the well-known fact that for $\G$ Karoubi-localising, $\N\G(\C)$ is also (a shift of) the fibre $\G(\mathrm{Nil}(\C)) \to \G(\C)$ where $\mathrm{Nil}(\C)$ is the category of nilpotent endomorphisms on $\C$, see e.g.\ \cite{LT2}*{Corollary 4.5 \& Prop.\ 4.7}.
In particular, if $\F\colon \Catp \to \A$ is a functor defined on Poincaré $\infty$-categories, we may form $\N\F^{\hyp}$, where $\F^\hyp = \F \circ \Hyp\colon \Catx \to \A$.
\end{construction}






\begin{definition}
\label{definition:assembly}%
The \defi{assembly map} for $\B\ZZ$ is the canonical colimit-exchange map associated to the constant diagram on $(\C,\QF)$ indexed on $\B\ZZ$:
\[
{\B\ZZ} \otimes \F(\C,\QF) := \colim\limits_{\B\ZZ} \F(\C,\QF) \xrightarrow{\ass} \F(\colim\limits_{\B\ZZ} (\C,\QF)) = \F((\C,\QF)_{\B\ZZ}).
\]
\end{definition}


\begin{remark}
\label{remark:circle-decompose-1}%
We always view $\B\ZZ$ as a pointed space.
The generator $1 \in \ZZ$ determines a pointed self homotopy $\tau$ of the composite $\ast_+ \to \ast \to \B\ZZ$, yielding a commutative square
\[
\begin{tikzcd}
*_+ \ar[r]\ar[d] & \ast \ar[d] \\
* \ar[r] & \B\ZZ 
\end{tikzcd}
\]
in pointed spaces, where we view $\tau$ as the commuting homotopy from the left-bottom composite to the top-right composite. This square is cocartesian, exhibiting $\B\ZZ$ as the suspension of $*_+$ in pointed spaces. 
For any stable $\infty$-category $\A$ and any object $x \in \A$, the above construction 
determines a decomposition
\[
\B\ZZ \otimes x = x \oplus \Sig x .
\]
In particular, if $(\C,\QF)$ is any Poincaré $\infty$-category then we can identify the assembly map with the map $\F(\C,\QF) \oplus \Sig\F(\C,\QF) \to \F((\C,\QF)_{\B\ZZ})$ induced by the resulting commutative square
\[
\begin{tikzcd}
\F(\C,\QF)\oplus \F(\C,\QF) \ar[r]\ar[d] & \F(\C,\QF) \ar[d] \\
\F(\C,\QF) \ar[r] & \F((\C,\QF)_{\B\ZZ}) 
\end{tikzcd}
\]
in $\A_{\F(\C,\QF)/}$, where we view the top left corner as living under $\F(\C,\QF)$ via the left summand.
\end{remark}




Let us now describe our circle decomposition theorem for Poincaré $\infty$-categories.
Given a Verdier-localising functor $\F\colon \Catp \to \A$ valued in a stable $\infty$-category $\A$ and a Poincaré $\infty$-category $(\C,\QF)$, 
consider the commutative diagram
\[
\begin{tikzcd}
\F(\Hyp(\C)) \ar[rr]\ar[dd]\ar[dr] && \F(\Hyp(\C_{\B\NN}))\ar[dd, dashed]\ar[dr] &\\
& \B\ZZ \otimes \F(\C,\QF) \ar[rr, "\ass", pos=0.4]\ar[dd] && \F((\C,\QF)_{\B\ZZ})\ar[dd]\ar[r] & \F(\C,\QF) \ar[dd] \\ 
\F(\Hyp(\C^{\natural})) \ar[rr, dashed]\ar[dr] && \F(\Hyp((\C^{\natural})_{\B\NN}))\ar[dr] &\\
& \B\ZZ \otimes \F(\C^{\natural},\QF) \ar[rr, "\ass",pos=0.4] && \F((\C^{\natural},\QF)_{\B\ZZ}) \ar[r] & \F(\C^{\natural},\QF)
\end{tikzcd}
\]
where the right diagonal arrows are instances of the map induced by the composite %
\[
\Hyp(\D_{\B\NN}) \to \Hyp(\D_{\B\ZZ}) \xrightarrow{\hyp} (\D,\Phi)_{\B\ZZ}
\]
and the left diagonal arrows are instances of the composite 
\[
\F(\Hyp(\D)) \xrightarrow{\hyp} \F(\D,\Phi) \xrightarrow{i} {\B\ZZ} \otimes \F(\D,\Phi).
\]
The main result of this section is then:


\begin{theorem}[The circle decomposition theorem]
\label{theorem:circle-decomposition}%
Let $\F\colon \Catp \to \A$ and $\G \colon \Catx \to \A$ be Verdier-localising functors valued in a stable $\infty$-category $\A$. %
Then the following hold:
\begin{enumerate}
\item
\label{item:main-face}%
The bottom face in the above diagram is cartesian. 
\item
\label{item:relative-circle}%
The functor $(\C,\QF) \mapsto \fib[\F((\C,\QF)_{\B\ZZ}) \to \F(\C,\QF)]$ on $\Catp$ is invariant under idempotent completion. In particular, 
the front right square in the above diagram is cartesian. 
\item
\label{item:nil}%
The functor $\D \mapsto \N\G(\D)$ on $\Catx$ is invariant under idempotent completion. In particular, the back face in the above diagram is cartesian. 
\end{enumerate}
In particular, for a Poincar\'e $\infty$-category $(\C,\QF)$ we have a natural decomposition
\[
\F((\C,\QF)_{\B\ZZ}) \simeq \F(\C,\QF) \oplus \Sig\F(\C^{\natural},\QF) \oplus \N\F^{\hyp}(\C).
\]
\end{theorem}


\begin{remark}
Unwinding definitions we find the map from right to left in the above natural decomposition is given by the sum of the following maps:
\begin{enumerate}
\item the tautological map $\F(\C,\QF) \rightarrow \F((\C,\QF)_{\B\ZZ})$,
\item the composite 
\begin{align*}
\Sig \F(\C^\natural,\QF) & \simeq \fib(\B\ZZ \otimes \F(\C^\natural,\QF) \rightarrow \F(\C^\natural,\QF)) \xrightarrow{\ass} \fib[\F((\C^\natural,\QF)_{\B\ZZ}) \rightarrow \F(\C^\natural,\QF)] \\
		& \simeq \fib[\F((\C,\QF)_{\B\ZZ}) \rightarrow \F(\C,\QF)] \longrightarrow \F((\C,\QF)_{\B\ZZ}),
\end{align*}
\item the map $\N\F^\hyp(\C) \rightarrow \F\Hyp(\C_{\B\NN}) \rightarrow \F((\C,\QF)_{\B\ZZ})$.
\end{enumerate}
Moreover, by construction the assembly map $\B\ZZ \otimes \F(\C,\QF) \to \F((\C,\QF)_{\B\ZZ})$ corresponds, under the splitting of its source as in 
Remark~\reffour{remark:circle-decompose-1},
to the composite of
\[
\F(\C,\QF) \oplus \Sig\F(\C,\QF) \to \F(\C,\QF) \oplus \Sig\F(\C^\natural,\QF) \to \F(\C,\QF) \oplus \Sig\F(\C^\natural,\QF) \oplus \N\F^{\hyp}(\C)
\]
where the first map is induced by the inclusion $(\C,\QF) \to (\C^\natural,\QF)$ and the second map is the summand inclusion. 
\end{remark}




\begin{corollary}[Ranicki-Shaneson decomposition]
\label{corollary:algebraic-shaneson-splitting}%
Let $(\C,\QF)$ be a Poincaré $\infty$-category and $\F\colon \Catp \to \A$ a bordism invariant Verdier-localising functor valued in a stable $\infty$-category $\A$. Then we have a natural decomposition
\[
\F((\C,\QF)_{\B\ZZ}) \simeq \F(\C,\QF) \oplus \Sig\F(\C^\natural,\QF)
\]
under which the assembly map from $\B\ZZ \otimes \F(\C,\QF) = \F(\C,\QF) \oplus \Sig\F(\C,\QF)$
is given by the identity on the first term and by the map induced by the inclusion $(\C,\QF) \to (\C^\natural,\QF)$ on the second term.
\end{corollary}

Arguing by induction we then conclude:
\begin{corollary}
\label{corollary:iterated-shaneson-splitting}%
Let $(\C,\QF)$ be a Poincaré $\infty$-category, $\F\colon \Catp \to \A$ a bordism invariant Verdier-localising functor valued in a stable $\infty$-category $\A$ and $n \geq 1$ an integer. Assume either that $\C_{\B\ZZ^{n-1}}$ is idempotent complete or that $\F$ is Karoubi-localising. Then the torus
assembly map
\[
\B\ZZ^n \otimes \F(\C,\QF) \to \F((\C,\QF)_{\B\ZZ^n})
\]
is an equivalence.
\end{corollary}



\begin{example}
\label{example:shaneson-splitting-Ghyp}%
If $\F\colon \Catp \to \A$ is a Verdier-localising functor we may form its genuine hyperbolisation $\F^\ghyp\colon \Catp \to \A^\gCt$, the genuine $\Ct$-$\A$-object valued functor whose underlying functor is $\F^\hyp = \F \circ \Hyp \colon \Catp \to \A^{\B\Ct}$, whose genuine fixed points is $\F$ and whose geometric fixed points is $\F^\bord$; see \S\reftwo{subsection:Mackey} for details. From loc.\ cit.\ we find $(\F^\ghyp)^\hyp \simeq \Ct \otimes \F^\hyp$ as functors $\Catx \to \A^{\gCt}$ where $C_2 \otimes - $ denotes the genuine induction functor, that is, the left adjoint of the underlying object functor $\A^{\gCt} \to \A$. Consequently,
Theorem~\reffour{theorem:circle-decomposition} gives an equivalence
\[
\F^{\ghyp}((\C,\QF)_{\B\ZZ}) \simeq \F^{\ghyp}(\C,\QF) \oplus \Sigma \F^{\ghyp}(\C^\natural,\QF) \oplus (\Ct \otimes \mathrm \N\F^\hyp(\C)).
\]
For $\F= \GW$, we have $\F^{\ghyp}= \KR$ and thus
\[
\KR((\C,\QF)_{\B\ZZ}) \simeq \KR(\C,\QF) \oplus \Sigma \KR(\C^{\natural},\QF) \oplus (\Ct \otimes \N\K(\C)),
\]
a Bass-Heller-Swan type decomposition for the real K-spectra of Poincaré categories. On underlying spectra with $\Ct$-action, this gives
\[
\K((\C,\QF)_{\B\ZZ}) \simeq \K(\C,\QF) \oplus \Sig \K(\C^\natural,\QF) \oplus [\Ct \otimes \N\K(\C)]
\]
with $\Ct$-action summand-wise. In particular, the $\Ct$-action swaps the two nil terms appearing in the classical Bass-Heller-Swan decomposition. On fixed points, we obtain
\[
\GW((\C,\QF)_{\B\ZZ}) = \GW(\C,\QF) \oplus \Sigma \GW(\C^{\natural},\QF) \oplus \N\K(\C)
\]
and on geometric fixed points, we obtain
\[
\L((\C,\QF)_{\B\ZZ}) = \L(\C,\QF) \oplus \Sigma \L(\C^{\natural},\QF).
\]
The same holds true for $\F=\KGW$, in which case the terms $\KK$ and $\KL$ appear instead of $\K$ and $\L$.
One may then insert the Poincar\'e $\infty$-categories appearing in Example~\reffour{example:tensor-BZ} and obtain formulas for the Grothendieck-Witt and L-theory of Laurent polynomial rings $R[t^{\pm 1}]$ equipped with involution sending $t$ to $t^{-1}$ or form parameter $\lambda[\ZZ]$ induced from a form parameter on $R$, as well as for the visible Poincar\'e structure on parametrised spectra over $X \times S^1$ with spherical fibration factoring through the projection $X \times S^1 \to X \to \Pic(\SS)$, see the circle decomposition part of the introduction.
\end{example}





\begin{remark}
\label{remark:all-involutions}%
Given a commutative ring $R$ with the identity involution, Example~\reffour{example:shaneson-splitting-Ghyp} gives in particular formulas for the Grothendieck--Witt- and L-spectra associated to the ring with $R[t^{\pm 1}]$ with involution given by $t\mapsto t^{-1}$. This is only one of the four canonical involutions $t \mapsto \pm t^{\pm 1}$. A number of things are known in the other cases as well:
\begin{enumerate}
\item
\label{item:twisted}%
The case $t\mapsto -t^{-1}$ was studied in~\cite{Milgram-Ranicki}, 
and fits in the more general setting of the twisted circle decomposition obtained in~\cite{twisted-shaneson}*{Theorem 1.5}, a special case of which yields the formulas
\[
\KGW^{\sym}(R[t^{\pm 1}];t\mapsto -t^{-1}) = \coeq[\KGW^{\sym}(R) \rightrightarrows \KGW^{\sym}(R)] \oplus \N\KK(R)
\]
and 
\[
\KL^{\sym}(R[t^{\pm 1}],t\mapsto -t^{-1}) = \coeq[\KL^{\sym}(R) \rightrightarrows \KL^{\sym}(R)] ,
\]
where one may, similarly to our circle decomposition theorem, replace the symmetric Poincaré structure
on $\Dperf(R)$ with the one determined by a form parameter on $R$, and the Poincaré structure on $\Dperf(R[t^{\pm 1}])$ with a corresponding induced form parameter on $R[t^{\pm1}]$.
The coequalisers in the above formulas are taken along the automorphisms $\id,\tau$ of $\KGW^{\sym}(R)$ and $\KL^{\sym}(R)$, respectively, where $\tau$ is the automorphism induced by the sign involution of $\QF^{\sym}_R$; in the case of $\KL^{\sym}(R)$ this is just the sign involution of $\KL^{\sym}(R)$, so that the coequaliser in question becomes the cofibre of multiplication by $2$ (compare~\cite{Milgram-Ranicki}). This coequaliser term can again be described as the image of a the assembly map, this time for a non-constant family of Poincaré $\infty$-categories indexed by the circle $\B\ZZ$.
\item
\label{item:motivic}%
The case where the involution is the identity can be deduced from the projective line formula of~\cite{motives}*{\S 6.1}. One obtains the formulas
\[
\KGW^{\sym}(R[t^{\pm 1}]) = \KGW^{\sym}(R) \oplus \Sig\KGW^{\sym,\scalebox{0.5}{[-1]}}(R) \oplus \N\KGW^{\sym}(R)\oplus \N\KGW^{\sym}(R)
\]
and 
\[
 \KL^{\sym}(R[t^{\pm 1}]) = \KL^{\sym}(R) \oplus \KL^{\sym}(R) \oplus \N\KL^{\sym}(R)\oplus \N\KL^{\sym}(R),
\]
where $\KGW^{\sym,\scalebox{0.5}{[-1]}}(R):=\KGW\big(\Dperf(R),(\QF^{\sym}_R)\qshift{-1}\big)$ and $\N\KGW^{\sym}(R)$ and $\N\KL^{\sym}(R)$ are the cofibres of the map $\KGW^{\sym}(R) \to \KGW^{\sym}(R[t])$ and $\KL^{\sym}(R) \to \KL^{\sym}(R[t])$, respectively. The situation becomes more transparent when passing to real $\K$-theory, in which case the above formulas can be combined into a single formula %
\[
\KKR^{\sym}(R[t^{\pm1}]) = [S^{\sig} \otimes \KKR^{\sym}(R)] \oplus \N\KKR^{\sym}(R) \oplus \N\KKR^{\sym}(R),
\]
where the assembly map term is replaced with tensoring with the sign representation sphere $S^{\sig}$, viewed as a genuine $\Ct$-space.
By \cite{motives}*{Theorem 6.3.1} the above nil groups vanish 
for $R$ regular noetherian of finite Krull dimension.
The analogous nil groups for quadratic $\L$-theory  %
have been shown to agree with Cappell's unitary nil terms, and are non-zero even for $R=\ZZ$, see \cite{Connolly-Ranicki}. 
\item One can expect to obtain a similar formula for the final involution $t\mapsto -t$, which can most likely be extracted from a suitable projective line formula in the setting of schemes with involution. This is the topic of work in progress~\cite{schemes-with-involution}.
\end{enumerate}
We remark that in the upcoming PhD thesis of Andreas Huber, a general formula for the value of Karoubi-localising invariants on certain pushouts of Poincaré $\infty$-categories is studied. 
One can then show that at least the involutions from~\reffouritem{item:twisted} and~\reffouritem{item:motivic}
give rise to Poincar\'e $\infty$-categories which can be described as pushouts of easier pieces referring only to $R$. In these cases, the results in Huber's thesis
apply to reproduce the above formulas in a unified manner. 
\end{remark}













\subsection{Proof of the circle decomposition theorem}
\label{subsection:proof}%
This subsection is dedicated to the proof of Theorem~\reffour{theorem:circle-decomposition}.
We shall use the following auxiliary construction.
Let \(\I\) be the category with objects \(\{-1,1\}\) and which has exactly two non-identity arrows \(\alp^+,\alp^-\), both going from \(-1\) to \(1\) (in particular, \(1\) and \(-1\) have no non-identity endomorphisms and there are no arrows from \(1\) to \(-1\)). 
The functor
\[
p\colon \I \to \B \ZZ
\]
sending $\alp^+$ to $1 \in \ZZ = \Hom_{\B\ZZ}(\ast,\ast)$ and $\alp^-$ to $0$ then exhibits $\B\ZZ$ as a localisation of $\I$.
\begin{remark}
\label{remark:unique}%
The functor $p$ above is completely determined by $p(1) = \ast$ and $p(\alp^+) \circ p(\alp^{-})^{-1} = 1 \in \ZZ$: More precisely the subgroupoid of $\Fun_*((\I,1),(\B\ZZ,\ast))$ spanned by those pointed functors $(\I,1) \to (\B\ZZ,\ast)$ such that $p(\alp^+) \circ p(\alp^-)^{-1} = 1$ is contractible. Indeed, this follows from the universal property of localisations and the fact that the group of pointed orientation preserving self-homotopy equivalences of the circle is contractible.

\end{remark}


Given a Poincaré $\infty$-category $(\C,\QF)$ we may tensor it with $\I$ in $\Cath$ to obtain a hermitian $\infty$-category $(\C,\QF)_\I$. Since the category \(\I\) is strongly finite,
Lemma~\refone{lemma:finite} and Proposition~\refone{proposition:tensor-strongly-finite} give that we may identify \(\C_\I\) with the stable \(\infty\)-category of diagrams \(\I\op \to \C\), such that the hermitian functor \(\QF_\I\) is given by
\[
\QF_{\I}(\vphi) = \colim_{i \in \I}\QF(\vphi(i)) 
\]
and its bilinear part of $\QF_{\I}$ is represented by the %
duality
\[
\Dual_{\I}(\vphi)(j) = \colim_{i \in \I}\Dual(\vphi(i))^{\Hom_{\I}(i,j)}.
\]
Calculating this explicitly gives that $\QF_\I$ takes the coequaliser of the two maps obtained by applying $\QF$, and 
\begin{equation}
\label{equation:explicit-dual-I}%
\Dual_{\I}
\big[
\begin{tikzcd}
\x  & \ar[l,shift right]\ar[l,shift left] \y 
\end{tikzcd}
\Big] = 
\Big[
\begin{tikzcd}
\Sig\Dual\x  & \ar[l,shift right]\ar[l,shift left] \Sig[\Dual\x \times_{\Dual \y}\Dual\x] 
\end{tikzcd}
\big],
\end{equation}
where the pullback $\Dual \x\times_{\Dual \y} \Dual \x$ is to be taken along the two evident maps, and the two arrows in the second diagram are given by the two projections from the fibre product. Direct inspection shows that this duality is perfect, so that $(\C,\QF)_\I$ is Poincaré.
We will also deduce this fact from abstract principles in Proposition~\reffour{proposition:I-tensor-poincare} below.




Let $\NN^{\pm} \subseteq \ZZ$ be the submonoid spanned by $\pm 1$, respectively. Let $p^{\pm}\colon \I \to \B\NN^{\pm}$ be the functor sending $\alp^{\pm}$ to $\pm 1 \in \NN^{\pm}$ and $\alp^{\mp}$ to $0$. Then the composites 
\[
\I \xrightarrow{p^{\pm}} \B\NN^{\pm} \to \B\ZZ
\]
both send the formal difference $\alp^+ \circ (\alp^-)^{-1}$ to $1 \in \ZZ$, and are hence both equivalent to $p\colon \I \to \B\ZZ$ by its uniqueness, see Remark~\reffour{remark:unique}. %
We thus obtain a commutative square of $\infty$-categories
\[
\begin{tikzcd}
\I \ar[r, "p^+"]\ar[d, "p^-"'] & \B \NN^+ \ar[d] \\
\B\NN^- \ar[r] & \B\ZZ
\end{tikzcd}
\]
in which all arrows are localisation functors: the horizontal functors exhibit their targets as localisations by $\alp^-$ and $-1$, respectively, while the vertical functors exhibit their targets as localisations by $\alp^+$ and $1$, respectively. We note that unlike $(\C,\QF)_\I$, the tensored hermitian $\infty$-categories $(\C,\QF)_{\B\NN^{\pm}}$ are generally not Poincaré. However, by Proposition~\reftwo{proposition:localizationforp4}\;\reftwoitem{item:induced-functor-verdier}, the functor $(\C,\QF)_\I \to (\C,\QF)_{\B\ZZ}$ is a Poincar\'e-Verdier projection, so we may use this projection, together with an analysis of its kernel, to infer information from $\F((\C,\QF)_\I)$ to $\F((\C,\QF)_{\B\ZZ})$.











The study of $(\C,\QF)_\I$ itself will in turn rest on a portion of %
the dual $\Q$-construction which we now recall. We consider the hermitian tensor $\dualQ_1(\C,\QF) = (\C,\QF)_{\Twar(\Del^1)}$ of $(\C,\QF)$ by the twisted arrow category of $\Del^1$. Since $\Twar(\Del^1)$ is the category classifying pushouts, for ease of notation we write $\mathcal{P} = \Twar(\Del^1)$ and denote its initial object by $-1$ and the two other objects by $1$ and $1'$.
Then $\dualQ_1(\C,\QF) = (\C,\QF)_{\mathcal{P}}$ is Poincaré by Lemma~\reftwo{lemma:dqexists} and similarly as with the tensor with $\I$, its objects are given by $\cP\op$-indexed diagrams $[x \rightarrow z \leftarrow y]$ with the Poincaré structure given by
\[
\QF_{\mathcal{P}}([x \rightarrow z \leftarrow y]) = \QF(x) \oplus_{\QF(z)} \QF(y)
\]
and duality given by 
\[
\Dual_{\mathcal{P}}([x \xrightarrow{f} z \xleftarrow{g} y]) = [\cof(\Dual(f)) \rightarrow \Sig \Dual(z) \leftarrow \cof(\Dual(g))].
\]
In fact, $\dualQ_1(\C,\QF)$ is part of a cosimplicial object $[n] \mapsto \dualQ_n(\C,\QF)$ in $\Catp$ with $\dualQ_0(\C,\QF) = (\C,\QF)$, see Lemma~\reftwo{lemma:dqexists} and what follows. The two face maps $[0] \to [1]$ (with images $\{0\}$ and $\{1\}$ respectively) and the degeneracy map $s\colon [1] \to [0]$ then induce Poincaré functors
\[
l,r\colon (\C,\QF) \to (\C,\QF)_{\mathcal{P}} \quad\quad s\colon (\C,\QF)_{\mathcal{P}} \to (\C,\QF)
\]
which are explicitly given by $l(x) = [x \rightarrow 0 \leftarrow 0]$ and $r(y) = [0 \rightarrow 0 \leftarrow y]$ and $s[x \rightarrow z \leftarrow y] = x \times_z y$. We note that $s$ admits a fully-faithful left adjoint given by $x \mapsto [x = x = x]$, and is hence in particular a split Verdier projection, in fact a split Poincar\'e-Verdier projection as one checks readily. On the other hand, the face functors $l,r$ are both split Poincaré-Verdier inclusions, with left adjoints given by evaluation at $1,1' \in \mathcal{P}$ respectively. %



\begin{construction}
\label{construction:twar}%
Consider the commutative diagram of categories
\begin{equation}
\label{equation:pre-square-I-Q1}%
\begin{tikzcd}
\{\ast, \ast\} \ar[r,"{1,1'}"] \ar[d]  & \mathcal{P} \ar[d] & \ar[l,"-1"'] \ast \ar[d,equal] \\
\ast \ar[r,"1"] & \I & \ar[l,"-1"'] \ast 
\end{tikzcd}
\end{equation}
where the left square is the pushout square exhibiting $\I$ as obtained from $\mathcal{P}$ by collapsing the objects $1,1'$ in $\mathcal{P}$ to the object $1 \in \I$, and the right square is a pullback square. By direct inspection the map $\mathcal{P} \to \I$ is a cocartesian fibration and hence the opposite of the right square satisfies the Beck-Chevalley property for right Kan extensions along the vertical maps. For a given stable $\infty$-category $\D$ we hence obtain a commutative diagram
\begin{equation}
\label{equation:stable-square-I-Q1}%
\begin{tikzcd}
\D \oplus \D \ar[r,"l \oplus r"]\ar[d] & \D_{\mathcal{P}} \ar[d] \ar[r,"\ev^{\cP}_{-1}"] & \D \ar[d,equal] \\
\D \ar[r,"i_1"] & \D_\I \ar[r,"\ev^{\I}_{-1}"] & \D
\end{tikzcd}
\end{equation}
where the left square is induced on tensors by the left square in the previous diagram and the right square is obtained by doing the same to the right square in the previous diagram and then replacing the horizontal arrows by their left adjoints. 
\end{construction}

When identifying tensors with diagram categories as above, all functors in the left square of~\eqreffour{equation:stable-square-I-Q1} are induced by right Kan extension and the horizontal maps in the right square are induced by evaluation 
at $-1$ in $\mathcal{P}$ and $\I$, respectively. Concretely, for the left square, $(l\oplus r)(x,y) = [x \rightarrow 0 \leftarrow y]$ and $i_1(x) = [0 \leftleftarrows x]$. 
Since tensoring with $\D$ is a colimit preserving operation the left square is cocartesian in $\Cat$.

\begin{proposition}
\label{proposition:I-tensor-poincare}%
The rows in diagram~\eqreffour{equation:stable-square-I-Q1} are split Verdier sequences. In addition, if $(\C,\QF)$ is a Poincar\'e $\infty$-category, the hermitian $\infty$-category $(\C,\QF)_\I$ is Poincar\'e and diagram~~\eqreffour{equation:stable-square-I-Q1} for $\D=\C$ refines to a commutative diagram in $\Catp$ whose rows are split Poincar\'e-Verdier sequences
\begin{equation}
\label{equation:square-I-Q1}%
\begin{tikzcd}
	(\C,\QF) \oplus (\C,\QF) \ar[r,"l\oplus r"]\ar[d] & (\C,\QF)_{\mathcal{P}} \ar[d] \ar[r,"\ev^{\cP}_{-1}"] & (\C,\QF\qshift{1}) \ar[d,equal]\\
	(\C,\QF) \ar[r,"i_1"] & (\C,\QF)_\I \ar[r,"\ev^{\I}_{-1}"] & (\C,\QF\qshift{1})
\end{tikzcd}
\end{equation}
where on underlying hermitian $\infty$-categories the left square is obtained from the left square of~\eqreffour{equation:pre-square-I-Q1} by tensoring with $(\C,\QF)$. In addition, the left square is a pushout in $\Catp$.
In particular, any additive functor $\F\colon \Catp \to \A$ to a stable $\infty$-category $\A$ sends the left square %
to an exact square. 
\end{proposition}
\begin{proof}
By direct inspection the rows of~\eqreffour{equation:stable-square-I-Q1} are fibre sequences. In addition, the evaluation functors $\ev^{\cP}_{-1}$ and $\ev^{\I}_{-1}$ admit fully-faithful left and right adjoints given by left and right Kan extension along the inclusion of $\{-1\}$. We hence conclude that the rows of~\eqreffour{equation:stable-square-I-Q1} are split Verdier sequences.

Given a Poincaré $\infty$-category $(\C,\QF)$, the left square of~\eqreffour{equation:stable-square-I-Q1} for $\D=\C$ refines to the left square of~\eqreffour{equation:square-I-Q1} on the level of hermitian $\infty$-categories by performing the tensor there, yielding again a cocartesian square since tensoring with $(\C,\QF)$ is colimit preserving. Since the top horizontal and left vertical arrows in the left square of~\eqreffour{equation:square-I-Q1} are Poincaré functors between Poincaré $\infty$-categories and $\Catp$ is closed in $\Cath$ under colimits by Proposition~\refone{proposition:Catp-cocomplete}, the left square in~\eqreffour{equation:square-I-Q1} is contained in $\Catp$ and is a pushout square there. In particular, $(\C,\QF)_\I$ is Poincaré.  
Now the horizontal arrows in the left square are Poincaré-Verdier inclusions by Proposition~\reftwo{proposition:example-cosieve}.
Combining this with the fact that the left square is cocartesian we conclude that the induced map $(\ev^{\cP}_{-1})_!\QF_{\cP} \Rightarrow (\ev^{\I}_{-1})_!\QF_\I$ is an equivalence. 
Concretely, using the left adjoints $x \mapsto [0 \rightarrow x \leftarrow 0]$ and $x \mapsto [0 \leftleftarrows x]$ of $\ev^{\cP}_{-1}$ and $\ev^{\I}_{-1}$, respectively, this map can be identified with the map
\[
\colim\big[\begin{tikzcd}
0 & \QF(x) \ar[r]\ar[l] & 0 
\end{tikzcd}\big] \xrightarrow{\simeq}
\colim\big[\begin{tikzcd}
\QF(x) \ar[r,shift right]\ar[r,shift left] & 0 
\end{tikzcd}\big] 
\]
induced on colimits by $\cP \to \I$. We now finish the proof by identifying the left hand side with $\QF\qshift{1}$.
\end{proof}

\begin{remark}
\label{rem:variant}%
\label{remark:dmet}%
The split Poincaré-Verdier projection $\ev^{\cP}_{-1}\colon (\C,\QF)_{\mathcal{P}} \to (\C,\QF\qshift{1})$ of Proposition~\reffour{proposition:I-tensor-poincare} factors through the Poincaré-Verdier quotient $\dMet_1(\C,\QF)=(\C,\QF)_{\cP}/l(\C,\QF)$,
see Definition~\reftwo{definition:dualmet}. Concretely, we may identify objects of $\dMet_1(\C,\QF)$ with arrows $z \leftarrow y$, such that its Poincaré structure is given by $[z \leftarrow y]\mapsto \cof[\QF(z) \to \QF(y)]$. Under this identification the projection $(\C,\QF)_{\cP} \to \dMet_1(\C,\QF)$ is given by
\[
[x \rightarrow z \leftarrow y] \mapsto [z \leftarrow y] \quad\quad \colim\big[\QF(x) \leftarrow \QF(z) \rightarrow \QF(y)\big]  \to \cof[\QF(z) \to \QF(y)] .
\]
The resulting Poincaré functor $\dMet_1(\C,\QF) \to (\C,\QF\qshift{1})$ is then the canonical projection 
\[
[z \leftarrow y] \mapsto z \quad\quad \cof[\QF(z) \to \QF(y)] \mapsto \QF\qshift{1}(z)
\]
from the dual metabolic category. Alternatively, identifying $\cof[\QF(x) \to \QF(z)]$ with $\fib[\QF\qshift{1}(x) \to \QF\qshift{1}(z)]$ we may also identify $\dMet_1(\C,\QF)$ with $\Met(\C,\QF\qshift{1})$ and the above projection with $\met\colon \Met(\C,\QF\qshift{1}) \to (\C,\QF\qshift{1})$. Summarising the situation, Diagram~\eqreffour{equation:square-I-Q1} can be extended to a diagram
\[
\begin{tikzcd}
	(\C,\QF) \ar[r,"l"] \ar[d,"i_l"] & (\C,\QF)_{\mathcal{P}} \ar[r,"\ev"] \ar[d,equal] & \Met(\C,\QF\qshift{1}) \ar[d,"\met"] \\
	(\C,\QF) \oplus (\C,\QF) \ar[r,"l\oplus r"]\ar[d] & (\C,\QF)_{\mathcal{P}} \ar[d] \ar[r,"\ev^{\cP}_{-1}"] & (\C,\QF\qshift{1}) \ar[d,equal]\\
	(\C,\QF) \ar[r,"i_1"] & (\C,\QF)_\I \ar[r,"\ev^{\I}_{-1}"] & (\C,\QF\qshift{1})
\end{tikzcd}
\]
in which also the top horizontal sequence is a split Poincar\'e-Verdier sequence, where $i_l$ is the inclusion of the left summand. 
\end{remark}









Having analysed $\D_\I$ and $(\C,\QF)_\I$, we now address the kernel of the induced functor $\D_\I \to \D_{\B\ZZ}$, which is a Verdier projection by Proposition~\reftwo{proposition:localizationforp4}.


\begin{construction}
\label{construction:E}%
Let $\D$ be a stable $\infty$-category.
For \(i \in \I\) and \(\x \in \D\) we let \(R_{i,\x}\colon \I\op \to \D\) denote the right Kan extension of \(\x\) along the inclusion \(\{i\} \subseteq \I\op\). Explicitly, the diagrams \(R_{1,\x}\) and \(R_{-1,\x}\) can be depicted as
\[
\begin{tikzcd}[row sep=2pt]
R_{1,x} = \big[0 & \ar[l,shift left]\ar[l,shift right] \x \big] && \text{and} && R_{-1,x} = \big[\x & \ar[l,shift left]\ar[l,shift right] \x \oplus \x \big], \
\end{tikzcd}
\]
respectively. 
The maps \(\alp^{\pm}\colon -1 \to 1\) in $\I$ now induce maps \(\alp^{\pm}_{\ast}\colon R_{-1,\x} \to R_{1,\x}\), whose fibres \(z^+_{\x} := \fib[\alp^+_{\x}]\) and \(z^-_x := \fib[\alp^-_{\x}]\) are depicted by the diagrams
\[
\begin{tikzcd}[row sep = 2ex]
z^+_x = \big[\x & \ar[l,shift left,"0"]\ar[l,shift right,"\id"'] \x\big] && \text{and} && z^-_x = \big[ \x & \ar[l,shift left,"\id",pos=0.45]\ar[l,shift right,"0"',pos=0.45] \x\big] , \
\end{tikzcd}
\]
respectively.
We then write \(\cN^{\pm}(\D) \subseteq \D_\I\) for the smallest stable and retract-closed subcategory of $\D_\I$ containing the objects \(z^{\pm}_{\x}\) for every \(\x \in \D\), and \(\cN(\D) \subseteq \D_{\I}\) for the smallest stable retract-closed subcategory of \(\D_{\I}\) containing $\cN^+(\D)$ and $\cN^-(\D)$. The categories $\cN^{\pm}(\D)$ are equivalent to the categories of nilpotent endomorphisms in $\D$ which we have briefly mentioned in Construction~\reffour{construction:nil}. If $(\C,\QF)$ is a Poincar\'e $\infty$-category,
the explicit formula above for \(\Dual_{\I}\) shows that \(\Dual_{\I}(z^+_{\x}) \simeq \Sigma z^-_{\Dual\x}\) and \(\Dual_{\I}(z^-_{\x}) \simeq \Sigma z^+_{\Dual \x}\), and so:
\begin{enumerate}
\item
\(\cN(\C)\) is %
invariant under the duality of $(\C,\QF)_\I$.
\item
The induced duality of $\cN(\C)$ switches between $\cN^+(\C)$ and $\cN^-(\C)$, and in particular determines an equivalence $\cN^-(\C) \simeq (\cN^+(\C))\op$.
\end{enumerate}
We will then denote by $\cN(\C,\QF) = (\cN(\C),(\QF_\I)_{|\cN(\C)})$ the Poincaré $\infty$-category obtained by restricting the Poincaré structure $\QF_\I$ to $\cN(\C)$.
\end{construction}


\begin{remarks}\
\label{remark:cofree}%
\begin{enumerate}
	\item The functor $\D_{\mathcal{P}} \to \D_\I$ in \eqreffour{equation:stable-square-I-Q1} sends $[x = x \leftarrow 0]$ to $z^+_x$ and $[0 \rightarrow x = x]$ to $z^-_x$. Indeed, this can be seen by direct inspection, or by noting that in both cases the object in question is obtained as a fibre of a map between two right Kan extended diagrams, and using the fact that the above induced functor is realised via right Kan extension.
	\item Since $\alp^{\mp}$ maps to an equivalence by $p^{\pm}\colon \I \to \B\NN^{\pm}$ %
and the induced map on tensors $p^{\pm}_*\colon \D_\I \to \D_{\B\NN^{\pm}}$ is realised via right Kan extension, we find for all $x \in \D$ that 
\[
p^+_*(z_x^-) = 0 \in \D_{\B\NN^+} \quad\text{and}\quad p^-_*(z^+_x) = 0 \in \D_{\B\NN^-} \, .
\]
\end{enumerate}
\end{remarks}



\begin{lemma}
\label{lemma:nil-term-hyperbolic}%
For a stable $\infty$-category $\D$, the full subcategories $\cN^+(\D),\cN^-(\D) \subseteq \cN(\D)$ induce a fully faithful and dense inclusion $\cN^-(\D) \oplus \cN^+(\D) \to \cN(\D)$. Moreover, for a Poincar\'e $\infty$-category $(\C,\QF)$, under the above identification $\cN^+(\C) = (\cN^-(\C))\op$, the full subcategory $\cN^-(\C) \oplus \cN^+(\C)$ of $\cN(\C,\QF)$ with the restricted Poincaré structure is naturally equivalent to $\Hyp(\cN^-(\C))$. 
In particular, if $\C$ is idempotent complete then the induced map 
\[
\Hyp(\cN^-(\C)) \to \cN(\C,\QF)
\]
is an equivalence.
\end{lemma}
\begin{proof}
We first show that for all $x^\pm \in \cN^\pm(\C)$, we have $\hom(x^+,x^-) = 0 = \map(x^-,x^+)$ from which it follows that the functor $\cN^-(\D) \oplus \cN^+(\D) \to \cN(\D)$ is fully faithful. Since its image is dense by definition of $\cN(\D)$, the first part of the lemma follows. To see the claim about the mapping spectra, it suffices to show $\map(z_x^-,z_y^+) = 0 = \map(z_x^+,z_y^-)$, since the mapping functor preserves finite colimits in each variable and retracts of $0$ are $0$. From the pushout formula $\I = \{1\} \amalg_{\{1,1'\}} \mathcal{P}$, recall diagram~\eqreffour{equation:pre-square-I-Q1}, we obtain the following pullback square
\[
\begin{tikzcd}
	\map_{\D_\I}(z_x^+,z_y^-) \ar[r] \ar[d] & \map_{\D_{\mathcal{P}}}([x=x\xleftarrow{0} x],[y\xrightarrow{0} y = y]) \ar[d]\\
	\map_\D(x,y) \ar[r,"\Delta"] & \map_\D(x,y) \times \map_\D(x,y)
\end{tikzcd}
\]
Unravelling the definitions, we find that the right vertical map is an equivalence in the first component of the target and the zero map in the second component of the target. As a consequence, the pullback vanishes as claimed. The same argument also shows $\map(z_x^-,z_y^+) = 0$.


For $(\C,\QF)$ a Poincar\'e $\infty$-category, using that the duality of $\cN(\C,\QF)$ interchanges $\cN^-(\C)$ and $\cN^+(\C)$, we deduce from the previous result on mapping spectra that $\B_\I$ vanishes upon restriction to $\cN^{\pm}(\C)$. As a result, $\QF_\I$ restricts to a linear functor on $\cN^\pm(\C)$, so we deduce from the computation
\[
\QF_\I(z^{+}_{\x}) = \colim \Big[
\begin{tikzcd}
\QF(\x) \ar[r,shift right,"0"']\ar[r,shift left,"\id"] & \QF(\x) 
\end{tikzcd}
\Big] = 0,
\]
and similarly for \(z^-_{\x}\), that $\QF_\I$ vanishes upon restriction to $\cN^{\pm}(\C)$. 
Hence, the restriction of $\QF_\I$ on $\cN^-(\C) \oplus \cN^+(\C)$ vanishes on each of the components and has a duality which swaps them, so we obtain that the induced Poincaré functor 
\[
\Hyp(\cN^-(\C)) \to (\cN^-(\C) \oplus \cN^+(\C),(\QF_\I)_{|\cN^+(\C)\oplus \cN^-(\C)})
\]
is an equivalence as desired.
\end{proof}
 




\begin{proposition}
\label{proposition:PV-sequences-Q1-hyp}%
For a stable $\infty$-category $\D$ the induced functor $p^+_*\colon \D_\I \to \D_{\B\NN^+}$ fits into a commutative diagram
\begin{equation}
\label{equation:tw-BN-C}%
\begin{tikzcd}
\D \ar[rr,"{x \mapsto [x = x \leftarrow 0]}"]\ar[d,"x \mapsto z_x^-"'] && \D_{\mathcal{P}} \ar[r]\ar[d] & \D_{\Del^1} \ar[d] \\
\cN^-(\D) \ar[rr] && \D_\I \ar[r, "p^+_*"] & \D_{\B\NN^+}
\end{tikzcd}
\end{equation}
in which the rows are Verdier sequences (with the top a split one) and the right square is induced on tensors by the square
\begin{equation}
\label{equation:tw-BN}%
\begin{tikzcd}
\mathcal{P} \ar[r,"c"]\ar[d] & \Del^1\ar[d] \\
\I \ar[r] & \B\NN^+ 
\end{tikzcd}
\end{equation}
where $c(-1) = c(1') =-1$ and $c(1) = 1$ (which admits a unique splitting).
\end{proposition}
\begin{proof}
Since the horizontal arrows in~\eqreffour{equation:tw-BN} are localisation functors, Proposition~\reftwo{proposition:localizationforp4} gives that the horizontal arrows in the right square of~\eqreffour{equation:tw-BN-C} are Verdier projections whose kernels are the smallest stable subcategories closed under retracts containing $[x = x \leftarrow 0] \in \D_{\mathcal{P}}$ and $z^-_x \in \D_\I$ for every $x \in \D$, respectively. In the former case the collection of cospans in $\D$ of the form $[x = x \leftarrow 0]$ already form a reflective and coreflective stable subcategory of $\D_{\mathcal{P}}$ which is equivalent to $\D$ itself via the displayed exact functor $x \mapsto [x = x \leftarrow 0]$. In the latter case the relevant smallest full stable subcategory closed under retracts is $\cN^-(\D)$. The left most vertical leg in the commutative diagram is necessarily the induced map on fibres, which by Remark~\reffour{remark:cofree} is given by the exact functor $\D \to \cN^-(\D)$ sending $x$ to $z_x^-$.
\end{proof}



\begin{proposition}
\label{proposition:PV-sequences-Q1}%
For a Poincaré $\infty$-category $(\C,\QF)$ the induced hermitian functor $p_*\colon (\C,\QF)_\I \to (\C,\QF)_{\B\ZZ}$ is Poincaré and furthermore fits in a commutative diagram
\begin{equation}
\label{equation:tw-S1}%
\begin{tikzcd}
\Hyp(\C) \ar[r]\ar[d] & (\C,\QF)_{\mathcal{P}} \ar[r,"s"]\ar[d] & (\C,\QF) \ar[d] \\
\cN(\C,\QF) \ar[r] & (\C,\QF)_\I \ar[r,"p_*"] & (\C,\QF)_{\B\ZZ} 
\end{tikzcd}
\end{equation}
whose rows are Poincaré-Verdier sequences (with the top a split one), whose right square is induced on tensors by the square
\[
\begin{tikzcd}
\mathcal{P} \ar[r]\ar[d] & \ast \ar[d] \\
\I \ar[r] & \B\ZZ \ ,
\end{tikzcd}
\]
and such that the two Poincaré functors out of $\Hyp(\C)$ are induced by the exact functors $x \mapsto [x \rightarrow x \leftarrow 0]$ and $x \mapsto z^-_x$, respectively.
\end{proposition}
\begin{proof}
Since $\cN(\C) \subseteq \C_\I$ is closed under the duality of $(\C,\QF)_\I$ we have by Proposition~\reftwo{proposition:localizationforp4}\;\reftwoitem{item:induced-functor-verdier} that $(\C,\QF)_\I \to (\C,\QF)_{\B\ZZ}$ is a Poincaré-Verdier projection whose kernel coincides with $\cN(\C,\QF)$. This gives the bottom row of~\eqreffour{equation:tw-S1}.
For the top row, note that the hermitian functor induced on tensors by the terminal map $\mathcal{P} \to \ast$ is the degeneracy functor 
$s$, which we observed earlier is a split Poincaré-Verdier projection. 
Its kernel consists of cospans of the form
$[x \rightarrow x \oplus y \leftarrow y] = [x = x \leftarrow 0] \oplus [0 \rightarrow y = y]$.
Since the Poincaré structure of $(\C,\QF)_{\mathcal{P}}$ vanishes on cospans of the form $[x = x \leftarrow 0]$ and on cospans of the form $[0 \rightarrow y = y]$, and at the same time the duality $\Dual_{\mathcal{P}}$ interchanges these two types of cospans, we conclude that the kernel of the Verdier projection $s\colon \C_{\mathcal{P}} \to \C$ is closed under the duality and identifies with $\Hyp(\C)$ when equipped with the restricted Poincaré structure. Example~\reftwo{example:kan-poincare} then yields the top Poincaré-Verdier sequence in \eqreffour{equation:tw-S1} and left most vertical leg in \eqreffour{equation:tw-S1} is necessarily the induced map on fibres, and by Remark~\reffour{remark:cofree} 
it is induced by the exact functor $\C \to \cN(\C)$ which sends $x$ to $z_x^-$.
\end{proof}


\begin{remark}
\label{remark:equ}%
Composing the Poincaré-Verdier projection $(\C,\QF)_\I \to (\C,\QF)_{\B\ZZ}$ of Proposition~\reffour{proposition:PV-sequences-Q1} with the induced Poincaré functor $(\C,\QF)_{\B\ZZ} \to (\C,\QF)$ yields a Poincaré functor
\[
\equ\colon (\C,\QF)_\I \to (\C,\QF)
\]
whose underlying hermitian functor is induced on tensors by the terminal map $\I \to \ast$. The Poincaré functor $\equ$ then provides a retraction to the split Poincaré-Verdier inclusion $i_1\colon (\C,\QF) \to (\C,\QF)_\I$ of Proposition~\reffour{proposition:I-tensor-poincare}, whose cofibre is identified with $(\C,\QF\qshift{1})$ as part of that proposition. In particular, 
for any additive functor $\F\colon \Catp \to \A$ we obtain a decomposition 
\[
\F((\C,\QF)_\I) = \F(\C,\QF) \oplus \F(\C,\QF\qshift{1}) .
\]
\end{remark}

\begin{remark}
\label{remark:I-is-assembly}%
Given a Poincaré $\infty$-category $(\C,\QF)$, the left square of Proposition~\reffour{proposition:I-tensor-poincare} and the right square of Proposition~\reffour{proposition:PV-sequences-Q1} fit into the following commutative diagram.
\[
\begin{tikzcd}
(\C,\QF) \oplus (\C,\QF) \ar[d]\ar[r] & (\C,\QF)_{\mathcal{P}} \ar[d]\ar[r] & (\C,\QF) \ar[d] \\
(\C,\QF) \ar[r] & (\C,\QF)_\I \ar[r] & (\C,\QF)_{\B\ZZ}
\end{tikzcd}
\]
Applying a Verdier-localising functor $\F\colon \Catp \to \cA$ with a stable target to this diagram sends the left square to a cartesian square by Proposition~\reffour{proposition:I-tensor-poincare}, and hence induces an equivalence
\[
\begin{tikzcd}
\F(\C,\QF) \displaystyle\mathop{\oplus}_{\F(\C,\QF) \oplus \F(\C,\QF)} \F(\C,\QF)\ar[dr,"\ass"]\ar[rr,"\simeq"] && \F(\C,\QF) \displaystyle\mathop{\oplus}_{\F((\C,\QF)_{\mathcal{P}})} \F((\C,\QF)_\I)  \ar[dl] \\
& \F((\C,\QF)_{\B\ZZ}) &
\end{tikzcd}
\]
from the assembly map on the left to the diagonal map on the right (see Remark~\reffour{remark:circle-decompose-1}). If furthermore $\F$ is bordism invariant then the map $\F((\C,\QF)_{\mathcal{P}}) \to \F(\C,\QF)$ is an equivalence by Proposition~\reffour{proposition:PV-sequences-Q1} so that the assembly map reduces to the map induced by the Poincaré functor $(\C,\QF)_\I \to (\C,\QF)_{\B\ZZ}$.

 
We then note that the inclusion $i_l\colon (\C,\QF) \to (\C,\QF) \oplus (\C,\QF)$ of the left summand induces a split Poincar\'e-Verdier inclusion of $(\C,\QF)$ into all four terms of the left square in the first displayed diagram above. Taking the quotients of each of these split Poincar\'e-Verdier inclusions, using Remark~\reffour{rem:variant}, we obtain the metabolic Poincar\'e-Verdier sequence
\[
\begin{tikzcd}
	(\C,\QF) \ar[r] \ar[d] & \Met(\C,\QF\qshift{1}) \ar[d] \\
	0 \ar[r] & (\C,\QF\qshift{1})
\end{tikzcd}
\]
It follows that both terms of the horizontal equivalence in the above commutative triangle receive a map from $\F(\C,\QF)$ which is split (see Remark~\reffour{remark:equ}) and whose induced map between cofibres is the equivalence
\[
\begin{tikzcd} \Sigma \F(\C,\QF) \ar[r,"\simeq"] & \cof[\Met(\C,\QF\qshift{1}) \to (\C,\QF\qshift{1})] \end{tikzcd}
\]
induced by the metabolic fibre sequence.
\end{remark}





We are now ready to establish the first part of the circle decomposition theorem (Theorem~\reffour{theorem:circle-decomposition}\;\reffouritem{item:main-face}):
	
\begin{theorem}
\label{theorem:circle-decomposition-idem}%
Let $\F\colon \Catp \to \A$ a Verdier-localising functor valued in a stable $\infty$-category $\A$ and $(\C,\QF)$ an idempotent complete Poincaré $\infty$-category. Then the square
\[
\begin{tikzcd}
\F(\Hyp(\C)) \ar[r]\ar[d] & \F(\Hyp(\C_{\B\NN^+}))\ar[d] \\
{\B\ZZ} \otimes \F(\C,\QF) \ar[r, "\ass"] & \F((\C,\QF)_{\B\ZZ}),
\end{tikzcd}
\]
is cartesian.
\end{theorem}
\begin{proof}
Consider the following vertical transformation between two diagrams of Poincaré $\infty$-categories:
\[
\begin{tikzcd}
0 \ar[rr]\ar[dd]\ar[dr] && \Hyp(\C\oplus\C)\ar[dd,dashed] \ar[rr,equal]\ar[dr] && \Hyp(\C \oplus \C)\ar[dd,dashed] \ar[dr] & \\
& \Hyp(\C)\ar[rr]\ar[dd] && \Hyp(\C_{\mathcal{P}}) \ar[rr]\ar[dd] && \Hyp(\C_{\Del^1})\ar[dd] \\
0 \ar[rr, dashed]\ar[dr] && \Hyp(\C)  \ar[rr,equal,dashed]\ar[dr] && \Hyp(\C)\ar[dr]  & \\
& \Hyp(\cN^-(\C))\ar[rr] && \Hyp(\C_\I) \ar[rr] && \Hyp(\C_{\B\NN^+}) \\
&& \color{white}{\;} \ar[r,phantom,"\rotatebox{270}{$\Rightarrow$}"] & \color{white}{\;}  && \\
0 \ar[rr]\ar[dd]\ar[dr] && (\C,\QF) \oplus (\C,\QF)\ar[dd,dashed] \ar[rr,equal]\ar[dr] && (\C,\QF) \oplus (\C,\QF)\ar[dd,dashed] \ar[dr] & \\
& \Hyp(\C)\ar[rr]\ar[dd] && (\C,\QF)_{\mathcal{P}} \ar[rr]\ar[dd] && (\C,\QF)\ar[dd] \\
0 \ar[rr, dashed]\ar[dr] && (\C,\QF)  \ar[rr,equal,dashed]\ar[dr] && (\C,\QF)\ar[dr]  & \\
& \cN(\C,\QF)\ar[rr] && (\C,\QF)_\I \ar[rr] && (\C,\QF)_{\B\ZZ}.
\end{tikzcd}
\]
The front horizontal rows are the Poincaré-Verdier sequences of Propositions \reffour{proposition:PV-sequences-Q1-hyp} (after applying $\Hyp$) and \reffour{proposition:PV-sequences-Q1}, and the middle squares in the top and bottom diagrams are those of Construction~\reffour{construction:twar} (after applying $\Hyp$) and Proposition~\reffour{proposition:I-tensor-poincare}, respectively.
The functor $\F$ hence sends each of the rows to a fibre sequence. In addition, the transformation between the top and bottom diagrams %
is an equivalence on left most faces by Lemma~\reffour{lemma:nil-term-hyperbolic} (which is where we use the assumption that $\C$ is idempotent complete), while the middle faces are both sent to cartesian squares via $\F$ by Proposition~\reffour{proposition:I-tensor-poincare}. %
After applying $\F$, the transformation from the top diagram to the bottom one thus induces an equivalence on total cofibres of the right most faces. At the same time, the bottom right most face encodes the assembly map, see Remark~\reffour{remark:circle-decompose-1}.
We conclude that the resulting square
\[
\begin{tikzcd}
\F(\Hyp(\C)) \displaystyle\mathop{\oplus}_{\F(\Hyp(\C \oplus \C))} \F(\Hyp(\C_{\Del^1})) \ar[r]\ar[d] & \F(\Hyp(\C_{\B\NN^+}))\ar[d] \\
{\B\ZZ} \otimes \F(\C,\QF) \ar[r] & \F((\C,\QF)_{\B\ZZ}) 
\end{tikzcd}
\]
is cartesian.
To finish the proof it will now suffice to show that the map
\[
\F((\Hyp(\C \oplus \C)) \to \F(\Hyp(\C_{\Del^1})),
\]
induced by the boundary inclusion $\{\ast,\ast\} \to \Del^1$, is an equivalence.
Indeed, this map 
is obtained by applying the Verdier-localising functor $\F^{\hyp}$ to the middle vertical map in the commutative rectangle
\[
\begin{tikzcd}
\C \ar[rr, "{x\mapsto (x,0)} "]\ar[d,equal] && \C \oplus \C\ar[rr, "{(x,y) \mapsto y}"]\ar[d] && \C \ar[d,equal] \\
\C \ar[rr, "{x \mapsto [x \to 0]}"] && \C_{\Del^1} \ar[rr, "{[x\to y] \mapsto y}"] && \C
\end{tikzcd}
\]
whose rows are (split) Verdier sequences. 
\end{proof}



\begin{remark}
\label{remark:proof-for-bordism-invariant}%
The proof of Theorem~\reffour{theorem:circle-decomposition-idem} described above simplifies when $\F$ is bordism invariant: in this case, $\F$ vanishes on the top diagram
as well as on the left most face in the bottom diagram. Since the rows are Poincaré-Verdier sequences we obtain an equivalence between the value of $\F$ on the remaining faces:
\[
\begin{tikzcd}
\F(\C,\QF) \oplus \F(\C,\QF)\ar[dd] \ar[dr] && \F(\C,\QF) \oplus \F(\C,\QF)\ar[dd] \ar[dr] & \\
& \F((\C,\QF)_{\mathcal{P}}) \ar[dd]\ar[dr,phantom,"\simeq"] && \F(\C,\QF)\ar[dd] \\
\F(\C,\QF)  \ar[dr] && \F(\C,\QF)\ar[dr]  & \\
& \F((\C,\QF)_\I)  && \F((\C,\QF)_{\B\ZZ}).
\end{tikzcd}
\]
Since the left one is exact so is the right, and we obtain that the assembly map is an equivalence. 
\end{remark}






 










Next, we address the parts of the circle decomposition theorem pertaining to the behaviour under idempotent completion, that is Theorem~\reffour{theorem:circle-decomposition}\;\reffouritem{item:relative-circle} and~\reffouritem{item:nil}.

\begin{lemma}
\label{lemma:S1-natural}%
For a stable $\infty$-category $\D$ and a Poincar\'e $\infty$-category $(\C,\QF)$, the squares
\[
\begin{tikzcd}
\D_{\B\NN} \ar[r] \ar[d] & \D \ar[d] & (\C,\QF)_{\B\ZZ} \ar[r] \ar[d] & (\C,\QF) \ar[d] \\
(\D^{\natural})_{\B\NN} \ar[r] & \D^{\natural} & (\C^\natural,\QF)_{\B\ZZ} \ar[r] & (\C^\natural,\QF)
\end{tikzcd}
\]
are cartesian in $\Catx$ and $\Catp$, respectively. 
\end{lemma}
\begin{proof}
The vertical maps in the squares are dense Karoubi and Poincaré-Karoubi inclusions, respectively.
At the same time, the horizontal arrows induce bijections on $\K_0$: indeed, they admit sections induced on tensors by the base point inclusions $\ast \to \B\NN$ and $\ast \to \B\ZZ$, and these section are surjective on $\K_0$ since their essential image generates their respective targets under finite colimits and desuspensions.
The desired result hence follows from Thomason's classification of dense subcategories.
\end{proof}





\begin{construction}
\label{construction:PV-square}%
Let $\D$ be a stable $\infty$-category. We let $\D^{\flat}_{\I} \subseteq \D_\I^\natural = (\D^{\natural})_\I$ be the dense subcategory spanned by those objects whose equaliser lies in $\D$. Similarly, given a Poincaré $\infty$-category 
$(\C,\QF)$ we define $(\C,\QF)^{\flat}_\I$ to be 
the Poincaré subcategory of $(\C^\natural,\QF)_\I$ spanned by $\C^{\flat}_\I$. We note that this subcategory is invariant under the duality since taking equalisers is duality invariant, see Remark~\reffour{remark:equ}.
By Lemma~\reffour{lemma:S1-natural}, the pasting lemma, and the fact that Poincaré-Verdier projections are closed under base change, see Remark~\reftwo{remark:cocartesian}~\reftwoitem{item:PV-stable-by-pullback}, 
the Poincaré-Verdier projection $(\C,\QF)^{\flat}_{\I} \to (\C^\natural,\QF)_{\B\ZZ}$ of Proposition~\reffour{proposition:PV-sequences-Q1} and the Verdier projection $\D^{\flat}_\I \to (\D^{\natural})_{\B\NN}$ from Proposition~\reftwo{proposition:localizationforp4}\;\reftwoitem{item:induced-functor-verdier} extend to the following Verdier and Poincaré-Verdier squares, respectively.
\[
\begin{tikzcd}
\D^{\flat}_\I \ar[r, dashed]  \ar[d, hook] &  \D_{\B\NN}\ar[d,hook] & (\C,\QF)^{\flat}_\I \ar[r, dashed]  \ar[d, hook] & (\C,\QF)_{\B\ZZ} \ar[d,hook] \\
\D^\natural_\I \ar[r] & (\D^\natural)_{\B\NN} &  (\C^\natural,\QF)_\I \ar[r] & (\C^\natural,\QF)_{\B\ZZ}. 
\end{tikzcd}
\]
\end{construction}







\begin{lemma}
\label{lemma:PV-sequence-3}%
The Poincaré $\infty$-category $(\C,\QF)^{\flat}_\I$ sits in a commutative diagram
\[
\begin{tikzcd}
	(\C,\QF) \ar[r,"i_1"] \ar[d, hook] & (\C,\QF)^{\flat}_\I \ar[r,"\ev_{-1}"] \ar[l,bend right,"\equ"']  \ar[d,hook] & (\C^\natural,\QF\qshift{1}) \ar[d, equal] \\
	(\C^\natural,\QF) \ar[r,"i_1"]& (\C^\natural,\QF)_\I \ar[r,"\ev_{-1}"] \ar[l,bend right,"\equ"'] & (\C^\natural,\QF\qshift{1})
\end{tikzcd}
\]
whose top row is a Poincaré-Verdier sequence, and whose bottom row is the split Poincaré-Verdier sequence of
Proposition~\reffour{proposition:I-tensor-poincare}.
Similarly, if $\D$ is a stable $\infty$-category then we have a diagram in $\Catx$ with the analogous properties involving $\D^{\flat}_\I$ and $\D^{\natural}_\I$.
\end{lemma}
\begin{proof}
Since the vertical maps are Poincar\'e-Karoubi equivalences and the bottom row is a split Poincaré-Verdier sequence, we find that the top row is a Poincar\'e-Karoubi sequence. Since the top right horizontal functor is also essentially surjective and its kernel is precisely the image of $(\C,\QF)$ under $i_1$,
we conclude that the top row is a Poincaré-Verdier sequence. The same argument also applies on the level of stable $\infty$-categories.
\end{proof}

The following two corollaries finish the proof of Theorem~\reffour{theorem:circle-decomposition}.

\begin{corollary}
\label{corollary:S1-natural}%
Let $\F\colon \Catp \to \A$ be a Verdier-localising invariant with stable target $\A$ and $(\C,\QF)$ a Poincar\'e $\infty$-category. Then both squares in the diagram  
\[
\begin{tikzcd}
	(\C,\QF) \ar[r]\ar[d]  & (\C,\QF)_{\B\ZZ} \ar[d]\ar[r] & (\C,\QF) \ar[d] \\
	(\C^\natural,\QF) \ar[r] & (\C^\natural,\QF)_{\B\ZZ} \ar[r] & (\C^\natural,\QF)
\end{tikzcd}
\]
are sent to cartesian squares by $\F$. In particular, Theorem~\reffour{theorem:circle-decomposition}\;\reffouritem{item:relative-circle} holds true.
Similarly, if $\D$ is a stable $\infty$-category then the analogous diagram in $\Catx$ is sent to a diagram with cartesian squares by any Verdier-localising functor $\G\colon \Catx \to \A$.
\end{corollary}
\begin{proof}%



Since the horizontal composites in the diagram under investigation are the identity, it suffices to consider the left square, which in turn is obtained from pasting the following two squares:
\[
\begin{tikzcd}
	(\C,\QF) \ar[r,"i_1"] \ar[d] & (\C,\QF)^{\flat}_\I \ar[r] \ar[d] & (\C,\QF)_{\B\ZZ} \ar[d] \\
	(\C^\natural,\QF) \ar[r,"i_1"] & (\C^\natural,\QF)_\I \ar[r] & (\C^\natural,\QF)_{\B\ZZ}  
\end{tikzcd}
\]
Here, the left square 
is sent to a cartesian square by $\F$ as a consequence of Lemma~\reffour{lemma:PV-sequence-3}. The right square is a Poincaré-Verdier square, see Construction~\reffour{construction:PV-square}, and is hence sent to a cartesian square since $\F$ is Verdier-localising, giving the desired result. The same argument applies when $\D$ is a stable $\infty$-category and $\F$ is a Verdier-localising functor on $\Catx$.
\end{proof}


\begin{corollary}
\label{corollary:nil-is-invariant}%
Let $\G\colon \Catx \to \cA$ be a Verdier-localising functor valued in a stable $\infty$-category $\cA$ and $\D$ a stable $\infty$-category. Then both squares in the diagram
\[
\begin{tikzcd}
	\D \ar[r]\ar[d] & \D_{\B\NN} \ar[r] \ar[d] & \D \ar[d] \\
	\D^\natural \ar[r] & (\D^{\natural})_{\B\NN} \ar[r] & \D^\natural
\end{tikzcd}
\]
are sent to cartesian squares by $\G$. In particular, Theorem~\reffour{theorem:circle-decomposition}\;\reffouritem{item:nil} holds true.
\end{corollary}
\begin{proof}
As before, it suffices to consider the left square, which is again obtained from pasting the following two squares:
\[
\begin{tikzcd}
	\D \ar[r] \ar[d] & \D^\flat_\I \ar[r] \ar[d] & \D_{\B\NN} \ar[d] \\
	\D^\natural \ar[r] & \D^\natural_\I \ar[r] & (\D^\natural)_{\B\NN}
\end{tikzcd}
\]
The left square is sent to a cartesian square by $\G$ as a consequence of (the stable $\infty$-category version of) Lemma~\reffour{lemma:PV-sequence-3}. The right square is a Verdier square, see Construction~\reffour{construction:PV-square}, and is hence sent to a cartesian square since $\G$ is Verdier-localising, giving the desired result.
\end{proof}















\subsection{The extended assembly map}
\label{subsection:extended}%

When $\C$ is idempotent complete or $\F$ Karoubi-localising Theorem~\reffour{theorem:circle-decomposition} can be expressed as saying that the assembly map $\B\ZZ \otimes \F(\C,\QF) \to \F((\C,\QF)_{\B\ZZ})$ is a split inclusion with complementary summand given by the nil term $\N\F^{\hyp}(\C)$ and therefore an equivalence if $\F$ is in addition bordism invariant. 

When $\C$ is not idempotent complete one can still obtain such a formulation by replacing the assembly map with the base change  
\[
[\B\ZZ \otimes \F(\C^{\natural},\QF)] \times_{(\C^{\natural},\QF)_{\B\ZZ}} (\C,\QF)_{\B\ZZ} \to (\C,\QF)_{\B\ZZ}
\]
of the assembly map at $(\C^{\natural},\QF)$.
Our goal in this subsection is to give a $\Catp$-level description for the domain of this map, which we shall henceforth call the \defi{extended assembly map}.
We will make use of this description when %
analysing the difference between the algebraic and geometric decompositions in \S\reffour{subsection:splittings}.

\begin{construction}
\label{construction:ext-ass}%
Let $\F\colon \Catp \to \A$ be a direct sum preserving functor to a stable $\infty$-category $\A$. 
Denote by 
\[
(\C,\QF) \oplus^{\flat} (\C,\QF) \subseteq (\C^{\natural},\QF) \oplus (\C^{\natural},\QF)
\]
the full Poincaré subcategory spanned by those pairs $(x,y) \in \C^{\natural} \oplus \C^{\natural}$ such that $x \oplus y \in \C$, so that the sum map $(\C^\natural,\QF) \oplus (\C^{\natural},\QF) \to (\C^{\natural},\QF)$ restricts to a map $\oplus\colon(\C,\QF) \oplus^{\flat} (\C,\QF) \to (\C,\QF)$. We then define a new functor ${\B\ZZ} \otimes^{\flat}\F \colon \Catp \rightarrow \A$ via the pushout in the square
\[
\begin{tikzcd}
\F((\C,\QF) \oplus^{\flat} (\C,\QF)) \ar[r,"\F(\oplus)"]\ar[d,"\F(\oplus)"'] & \F(\C,\QF) \ar[d] \\
\F(\C,\QF) \ar[r] & {\B\ZZ} \otimes^{\flat} \F(\C,\QF) \ .
\end{tikzcd}
\]
Let us make explicit again that ${\B\ZZ} \otimes^{\flat} \F(\C,\QF)$ is not given by some operation applied to $\F(\C,\QF)$. We note that the inclusion $(\C,\QF) \oplus (\C,\QF) \to (\C,\QF) \oplus^\flat (\C,\QF)$ induces a map $\B\ZZ \otimes \F(\C,\QF) \to \B\ZZ \otimes^\flat \F(\C,\QF)$ which is an equivalence if $\C$ is idempotent complete.
The square
\[
\begin{tikzcd}
(\C,\QF) \oplus^{\flat} (\C,\QF) \ar[r]\ar[d] & (\C,\QF) \ar[d] \\
(\C,\QF) \ar[r] & (\C,\QF)_{\B\ZZ} \ ,
\end{tikzcd}
\]
obtained via restriction from the pushout square of $(\C^{\natural},\QF)_{\B\ZZ}$ obtained from writing $\B\ZZ$ as the pushout $\ast \amalg_{\{\ast,\ast\}} \ast$, then determines a map
\[
\ass^{\flat}\colon {\B\ZZ} \otimes^{\flat} \F(\C,\QF) \to \F((\C,\QF)_{\B\ZZ}) ,
\]
whose precomposition with the above mentioned tautological map $\B\ZZ \otimes \F(\C,\QF) \to \B\ZZ \otimes^\flat \F(\C,\QF)$ is the assembly map. 
\end{construction}

We %
call $\ass^{\flat}$ %
the extended assembly map. It coincides with base change of the assembly map of $(\C^{\natural},\QF)$ by the next remark.
If $\C$ is idempotent complete, the extended assembly map recovers the assembly map.



\begin{remark}
\label{remark:circle-decompose-2}%
We have a Poincaré-Verdier sequence
\[
(\C,\QF) \xrightarrow{x \mapsto (x,0)} (\C,\QF) \oplus^{\flat} (\C,\QF) \xrightarrow{(x,y) \mapsto y} (\C^{\natural},\QF)
\]
whose inclusion is retracted by the Poincaré functor $\oplus \colon (\C,\QF) \oplus^{\flat} (\C,\QF) \to (\C,\QF)$. We thus obtain natural decompositions
\[
\B\ZZ \otimes^{\flat} \F(\C,\QF) = \F(\C,\QF) \oplus \Sig\F(\C^{\natural},\QF) \quad\text{and}\quad \F((\C,\QF) \oplus^{\flat} (\C,\QF)) = \F(\C,\QF) \oplus \F(\C^{\natural},\QF)
\]
and conclude that the external rectable in %
\[
\begin{tikzcd}
\B\ZZ \otimes^{\flat} \F(\C,\QF) \ar[r]\ar[d]& \F((\C,\QF)_{\B\ZZ})\ar[r]\ar[d] & \F(\C,\QF) \ar[d] \\
\B\ZZ \otimes \F(\C^\natural,\QF) \ar[r] & \F((\C^{\natural},\QF)_{\B\ZZ})\ar[r]& \F(\C^\natural,\QF)
\end{tikzcd}
\]
is cartesian with horizontal fibres given by $\Sig\F(\C^{\natural},\QF)$. It then follows that both squares are cartesian by Corollary~\reffour{corollary:S1-natural}.


\end{remark}


Unravelling the definitions, we find that the circle decomposition theorem now implies: 
\begin{corollary}
For any Poincar\'e $\infty$-category $(\C,\QF)$ and any Verdier-localising functor $\F$, the extended assembly map 
\[
\B\ZZ \otimes^\flat \F(\C,\QF) \to \F((\C,\QF)_{\B\ZZ})
\]
is a split inclusion with complementary term the nil term $\N\F^\hyp(\C)$. In particular, the extended assembly map is an equivalence if $\F$ is bordism invariant.
\end{corollary}




\begin{remark}
\label{remark:algebraic-splitting-2}%
\label{remark:algebraic-splitting}%
Recall from 
Remark~\reffour{remark:I-is-assembly} that for $\F$ a Verdier-localising functor and $(\C,\QF)$ a Poincaré $\infty$-category 
the assembly map can be identified with the induced map
\[
\F(\C,\QF) \displaystyle\mathop{\oplus}_{\F((\C,\QF)_{\mathcal{P}})} \F((\C,\QF)_\I)  \to \F((\C,\QF)_{\B\ZZ}).
\]
Let us complement this by explaining how to do the same for the extended assembly map.
For this, let $(\C,\QF)^{\flat}_{\mathcal{P}}$ be the dense Poincaré subcategory of $(\C^\natural,\QF)_{\mathcal{P}}$ spanned by those objects $[x \rightarrow z \leftarrow y] \in \C^{\natural}_{\mathcal{P}}$ whose image $x \times_z y \in \C^{\natural}$ under the degeneracy functor $s$ belongs to $\C$. By construction, the diagram of Proposition~\reffour{proposition:I-tensor-poincare} and Remark~\reffour{rem:variant} then restricts to a commutative diagram
\[
\begin{tikzcd}
(\C,\QF) \ar[r] \ar[d,"i_r"] & (\C,\QF)^\flat_{\mathcal{P}} \ar[r,"\ev"] \ar[d,equal] & \Met(\C^\natural,\QF\qshift{1}) \ar[d] \\
(\C,\QF)\oplus^{\flat} (\C,\QF) \ar[d]\ar[r] & (\C,\QF)^{\flat}_{\mathcal{P}} \ar[d]\ar[r,"\ev_{-1}"] & (\C^{\natural},\QF\qshift{1}) \ar[d,equal] \\
(\C,\QF) \ar[r] & (\C,\QF)^{\flat}_\I \ar[r,"\ev_{-1}"] & (\C^{\natural},\QF\qshift{1})
\end{tikzcd}
\]
in which the rows are Poincaré-Verdier sequences by Lemma~\reffour{lemma:PV-sequence-3} (and the analogous argument for the top rows).
As in Remark~\reffour{remark:I-is-assembly}, we obtain a commutative triangle 
\[
\begin{tikzcd}
\F(\C,\QF) \displaystyle\mathop{\oplus}_{\F((\C,\QF) \oplus^{\flat} (\C,\QF))} \F(\C,\QF)\ar[dr,"\ass^{\flat}"]\ar[rr,"\simeq"] && \F(\C,\QF) \displaystyle\mathop{\oplus}_{\F((\C,\QF)^{\flat}_{\mathcal{P}})} \F((\C,\QF)^{\flat}_\I)  \ar[dl] \\
& \F((\C,\QF)_{\B\ZZ}) &
\end{tikzcd}
\]
whose horizontal functor induces an equivalence 
from the extended assembly map on the left to the diagonal map on the right. In particular, if $\F$ is bordism invariant then the extended assembly map is equivalent to the map induced by the Poincaré functor $(\C,\QF)^{\flat}_\I \to (\C,\QF)_{\B\ZZ}$. Again, just as in Remark~\reffour{remark:I-is-assembly}, it follows that splitting off the canonical summand $\F(\C,\QF)$ in the above horizontal equivalence, we obtain the canonical equivalence
\[
\Sigma\F(\C^\natural,\QF) \xrightarrow{\simeq} \cof[\F(\Met(\C^\natural,\QF\qshift{1})) \to \F(\C^\natural,\QF\qshift{1})]
\]
induced by the metabolic fibre sequence.
\end{remark}


\subsection{The framed transfer for the circle}
\label{subsection:splittings}%

In this section, we discuss an observation due to Ranicki that for L-spectra, in fact more generally for any bordism invariant Verdier-localising invariant $\F$, in addition to the Ranicki-Shaneson decomposition obtained in \reffour{corollary:algebraic-shaneson-splitting}, there is a second such decomposition which Ranicki refers to as the geometrically relevant decomposition. It is based on studying a suitable product with the visible signature of the circle $S^1 \simeq \B\ZZ$, thought of as a Poincar\'e element of $(\Spa^\free,(\QF^\uni)\qshift{-1})_{\B\ZZ}$ using a stable framing of $S^1$. Following Ranicki, we make explicit the difference between the two decompositions. It follows from this analysis that the decompositions agree when one uses the null-bordant stable framing and restricts attention to idempotent complete Poincar\'e $\infty$-category, but not in general.

Let us now explain the situation in more detail.
\begin{construction}
\label{const:transfer}%
To begin, recall that a finite connected Poincar\'e complex $X \in \Sps$, with dualising complex $\nu_X\colon X\to \Pic(\SS) \subset \Spa$, has a visible signature $\vsign(X)$, a Poincaré object in $(\Spa_X^\free,\QF^\vis_{\nu_X})$, see \S\refone{subsection:parametrised-spectra} or Appendix~\reffour{section:appendix-visible-revisited} below for a second perspective.
Given a Poincar\'e $\infty$-category $(\C,\QF)$, let us write 
\[
(\C, \QF)_{(X, \nu_X)} \coloneq \colim_{x\in X} (\C, \QF\otimes \nu_X(x)) = (\C,\QF) \otimes (\Spaf, \QF^\uni)_{(X, \nu_X)} = (\C,\QF) \otimes (\Spa_X^\free,\QF^\vis_{\nu_X})
\]
for the colimit of the canonical diagram induced by $\nu_X$. Recall that the tensor unit $(\Spaf,\QF^\uni)$ of $\Catp$ corepresents the Poincar\'e objects functor; In particular, $\vsign(X)$ is equivalently described by a Poincar\'e functor $(\Spaf,\QF^\uni) \to (\Spa_X^\free,\QF^\vis_{\nu_X})$. 
Tensoring the visible signature $\vsign(X)$ of $X$ with 
$(\C,\QF)$ then gives the Poincar\'e functor 
\[
(\C, \QF) \to (\C, \QF)\otimes (\Spaf, \QF^\uni)_{(X, \nu_X)} \simeq (\C, \QF)_{(X, \nu_X)},
\]
which we refer to as the \defi{parametrised transfer functor} associated to the Poincar\'e complex $X$. 
\end{construction}

\begin{remark}
\label{rem:framings}%
Suppose that the dualising complex $\nu_X$ is equipped with a framing, that is, an equivalence  %
\[
\tau\colon \nu_X \simeq \SS_X\qshift{-n}
\]
where $n=\dim(X)$ and $\SS_X\qshift{-n}$ is the constant local system on $X$ with value $\SS\qshift{-n} = \Omega^n\SS$.
Then we may form the composite, often referred to as \emph{spherical fundamental class} %
\[
c\colon \SS \longrightarrow \Th(\nu_X) \xrightarrow{\Th(\tau)} \Omega^n\Th(\SS_X) = \SS\qshift{-n}[X], 
\]
where for a stable spherical fibration $\xi\colon X\to \Pic(\SS)$, we write $\Th(\xi) = \colim_X \xi$ for its Thom spectrum and as above $\SS\qshift{-n}[X] = \Omega^n\SS[X]$. The first map in the above composite, often call the collapse map of $X$, is the composite of the constant inclusion $\SS\to \map(X, \SS)$ and the equivalence $\map(X, \SS)=\lim_X \SS\simeq \Th(\nu_X)$ issued from the defining property of the dualising complex. 


Conversely, any map $c\colon \SS \to \SS\qshift{-n}[X]$ which induces an equivalence
\[
\lim_X E \simeq \hom_X(\SS^{-n}_X,\SS^{-n}_X \otimes E) \xrightarrow{\colim} \hom(\SS^{-n}[X], \colim_X \SS_X^{-n} \otimes E) \xrightarrow{c} \colim_X \mathbb S^{-n}_X \otimes E
\]
for all $E \in \Sp_X$ determines a framing of $\nu_X$ by the uniqueness of the dualising complex; its associated spherical fundamental class then coincides with $c$. Passing between $\tau$ and $c$ in this manner, one obtains an equivalence between the $\infty$-groupoid framings of $\nu_X$ and the full subgroupoid of $\Omega^{\infty+n} \SS[X] = \Hom(\SS,\SS\qshift{-n}[X])$ spanned by the spherical fundamental classes. %
We also note that any spherical fundamental class $c$ induces, %
via the projection $q\colon X\to *$, an element 
\[
q_*[c]\in \pi_n(\SS)
\]
which we call the Pontryagin-Thom class of $c$.
If $X$ is a stably framed smooth manifold, then $q_*[c]$ is the element determined by the usual Pontryagin-Thom isomorphism between stable homotopy and framed bordism groups.  
\end{remark}

\begin{construction}
\label{construction:visible-sig-c}%
Given a spherical fundamental class $c\colon \SS \to \SS\qshift{-n}[X]$, we may view $X$ as a framed Poincar\'e complex as described in Remark~\reffour{rem:framings} and obtain an induced identification %
$(\C, \QF)_{(X, \nu_X)}\simeq (\C, \QF\qshift{-n})_X$ for any Poincaré $\infty$-category $(\C,\QF)$. We may hence view the visible signature of $X$ as a Poincar\'e object $\vsign_c(X)$ in $(\Sp^\free, (\QF^\uni)\qshift{-n})_X$ where we have added the subscript $c$ to keep track of the dependence of this construction on the spherical fundamental class $c$. 
By Proposition~\reffour{proposition:visible-signature-revisited},
under the embedding $\Spaf_X \subseteq \Fun(X\op,\Pro(\Spaf))$ (see Construction~\refone{construction:tensor} in the case of tensoring by an $\infty$-groupoid) the Poincaré object $\vsign_c(X)$ corresponds to the 
constant diagram $\SS_X$ equipped with the Poincar\'e form
\[
\beta_c\colon \SS \xrightarrow{\bet_{\uni}} \QF^{\uni}(\SS) \xrightarrow{c\otimes \QF^{\uni}(\SS)} \SS\qshift{-n}[X] \otimes \QF^{\uni}(\SS) = \colim_{X} \QF^{\uni}(\SS)\qshift{-n} = (\QF^\uni)_X\qshift{-n}(\SS_X) ,
\]
where $\bet_{\uni}$ is the universal Poincaré form on $\SS$. 
The unique map $q\colon X \to \ast$ induces a Poincar\'e functor
\[
q_*\colon (\Spaf,(\QF^{\uni})\qshift{-n})_{X} \longrightarrow (\Spaf,(\QF^{\uni})\qshift{-n})
\]
which we may apply to the visible signature to obtain a 
Poincaré object 
$q_*\vsign_c(X) = (q_*\SS_X,q_*\beta_c) = (\SS^{X},q_*\beta_c)$
in $(\Sp^\free,(\QF^\uni)\qshift{-n})$. 
Tensoring with $\vsign_c(X)$ and $q_*\vsign_c(X)$ then yields 
Poincar\'e functors
\[
\tr_{(X,c)}\colon (\C, \QF\qshift{n})\to (\C, \QF)_X \quad \text{and}\quad \tr_{(X,c)}^q \colon \quad (\C,\QF\qshift{n}) \to (\C,\QF)
\]
which we refer to as the \defi{(parametrised) framed transfer} associated to $X$ and its spherical fundamental class $c$. We note that the framed transfer is obtained from the parametrised framed transfer by post-composition with the map induced by $q\colon X \to \ast$.%
\end{construction}




We now apply these observations to $X=\B\ZZ$. Using the splitting 
\[
\SS\qshift{-1}[\B\ZZ] = \SS \oplus \SS\qshift{-1}
\]
of Remark~\reffour{remark:circle-decompose-1}
the $\infty$-groupoid of spherical fundamental classes identifies with the 
subgroupoid of $\Hom(\SS,\SS \oplus \SS\qshift{-1})$
spanned by those maps $c\colon \SS \to \SS \oplus \SS\qshift{-1}$ whose post-composition with the projection to $\SS$ is an equivalence. This $\infty$-groupoid therefore splits as a product $\gl_1(\SS) \times \Om^{\infty+1}\SS$, canonically pointed by the pair $(\id,0)$. The set of connected components of this $\infty$-groupoid is
\[
\pi_0\gl_1(\SS) \times \pi_1(\SS) = \{\pm1\} \times \ZZ/2 ,
\]
so that the equivalence type of any fundamental class is encoded by a sign 
$\eps(c) \in \{\pm1\}$,
which encodes the gap between the orientation determined by $c$ and our fixed orientation of $\B\ZZ$, and the associated Pontryagin-Thom class $q_*[c] \in \pi_1(\SS) = \ZZ/2$. 


\begin{remark}
\label{remark:pontryagin}%
Geometrically, the Pontryagin-Thom class determines whether or not 
the corresponding stable framing of the 1-circle extends to a stable framing of the 2-disc. 
From the present perspective, this is equivalent to the question whether or not the the map
$\SS \to \SS^{\B\ZZ}$ induced by $q\colon \B\ZZ \to \ast$ refines to a Lagrangian of the Poincaré object $q_*\vsign_c(\B\ZZ)$ in $(\Spaf,(\QF^{\uni})\qshift{-1})$.
Indeed, unwinding the definitions we get that the restriction of the form $q_*\beta_c \in (\QF^{\uni})\qshift{-1}(\SS^{\B\ZZ})$ %
along the constant inclusion $\SS\to \SS^{\B\ZZ}$ %
is given by the composite
\[
\SS \xrightarrow{q_*[c]} \SS^{-1} \to (\QF^{\uni})\qshift{-1}(\SS)
\]
where the second map is given by the universal Poincar\'e form. This composite is then null-homotopic if and only if the first map is, since $\SS \to \QF^{\uni}(\SS)$ admits a retraction $\QF^{\uni}(\SS) \to  \Lam^{\uni}(\SS) = \SS$ by projecting to the linear part of the universal Poincaré structure. In addition, when this composite is null-homotopic, any null-homotopy exhibits $\SS$ as a Lagrangian in $q_*\vsign_c(\B\ZZ)$, just like any extension of a stable framing of the 1-circle to the 2-disc constitutes a stably framed null-cobordism. In particular, if $q_*[c] = 0$, then $q_*\vsign_c(\B\ZZ)$ is null-cobordant. Conversely, it can also be shown that if $q_*[c] \neq 0$, then $q_*\vsign_c(\B\ZZ)$ is not null-cobordant.
\end{remark}




A specific feature of $\B\ZZ$ is that the $\K$-class underlying the visible signature $\vsign_c(\B\ZZ)$ is trivial. As a consequence, the (parametrised) framed transfers on $(\C,\QF)$ and $(\C^\natural,\QF)$, respectively, factor as follows:
\[
\begin{tikzcd}[column sep=large]
	(\C,\QF\qshift{1}) \ar[r] & (\C^\natural, \QF\qshift{1}) \ar[r,"{\tr^\flat_{(\B\ZZ,c)}}", dashed] & (\C, \QF)_{\B\ZZ} \ar[r,hook]  & (\C^\natural, \QF)_{\B\ZZ}
\end{tikzcd}
\]
We call $\tr^\flat_{(\B\ZZ,c)}$ the \defi{extended parametrised framed transfer} associated to $(\B\ZZ,c)$ and its composite $\tr^{\flat,q}_{(\B\ZZ,c)}$ with the tautological map $(\C,\QF)_{\B\ZZ} \to (\C,\QF)$ induced by $q \colon \B\ZZ \to \ast$ the \defi{extended framed transfer}.


\begin{proposition}
\label{proposition:geometric}%
Let $\F\colon \Catp \to \A$ be a Verdier-localising functor valued in a stable $\infty$-category $\A$. Under the equivalence of Theorem~\reffour{theorem:circle-decomposition}, the extended parametrised framed transfer
\[
\F(\C^\natural,\QF\qshift{1}) \xrightarrow{\tr^{\flat}_{(\B\ZZ,c)}} \F((\C, \QF)_{\B\ZZ})\simeq \F(\C,\QF) \oplus \Sig\F(\C^\natural,\QF) \oplus \N\F^{\hyp}(\C)
\]
has the following components:
\begin{enumerate}
\item
\label{item:base-component}%
Its component at $\F(\C,\QF)$ is induced by the extended framed transfer $\tr_{(\B\ZZ,c)}^{\flat,q}$ associated to $(\B\ZZ,c)$.
\item
\label{item:nil-component}%
Its component at $\N\F^{\hyp}(\C)$ is zero.
\item
\label{item:sig-component}%
Its component at $\Sig\F(\C^\natural,\QF)$ is given by $\eps(c)\cdot \partial$, where $\partial$ is the boundary map of the metabolic sequence. 
\end{enumerate}
\end{proposition}



\begin{corollary}%
\label{corollary:geometric-shaneson-splitting}%
Let $\F\colon \Catp \to \A$ be a bordism invariant and Verdier-localising functor valued in a stable $\infty$-category $\A$. 
Then the Poincaré functor $(\C,\QF) \to (\C,\QF)_{\B\ZZ}$ induced by 
the base point $\ast \to \B\ZZ$ and the extended framed transfer $\tr_{\B\ZZ}^{\flat,q}$ of $(\B\ZZ,c)$ induce an equivalence
\[
\F(\C,\QF) \oplus \Sig\F(\C^\natural,\QF) \simeq \F(\C,\QF) \oplus \F(\C^\natural,\QF\qshift{1}) \xrightarrow{\simeq} \F(\C,\QF)_{\B\ZZ} .
\]
\end{corollary}

We warn the reader again that this \defi{geometric Ranicki-Shaneson decomposition} is generally different from the one appearing in Corollary~\reffour{corollary:algebraic-shaneson-splitting}, as will discuss below.

\begin{proof}[Proof of Proposition~\reffour{proposition:geometric}]
Statement~\reffouritem{item:base-component} holds by definition. Next, we recall 
the Poincaré-Verdier projection
\[
p_*\colon (\Spaf,(\QF^{\uni})\qshift{-1})_{\I}^\flat \to (\Spaf,(\QF^{\uni})\qshift{-1})_{\B\ZZ}.
\]
We will show momentarily that the visible signature $\vsign_c(\B\ZZ)$ lifts 
to a Poincar\'e object $\bar{\sigma}^\vis_c(\I)$ in $(\Spaf,(\QF^{\uni})\qshift{-1})_{\I}^\flat$. As a consequence, the refined parametrised transfer factors as a composite
\[
(\C^\natural,\QF\qshift{1}) \to (\C,\QF)^\flat_\I \to (\C,\QF)_{\B\ZZ}.
\]
By Remark~\reffour{remark:algebraic-splitting-2}, we find that second displayed functor, after applying $\F$, factors as
\[
\F((\C,\QF)^\flat_\I) \to \B\ZZ\otimes^{\flat}\F(\C,\QF) \to \F((\C,\QF)_{\B\ZZ})
\]
where the second map is the extended assembly map of \S\reffour{subsection:extended}.
Since $\N\F^\hyp(\C)$ is the cofibre of the extended assembly map, we deduce statement~\reffouritem{item:nil-component}. 
For~\reffouritem{item:sig-component}, note that, by the last part of Remark~\reffour{remark:algebraic-splitting-2}, the left map in the last factorisation induces, after splitting off the summand $\F(\C,\QF)$ on both sides, the boundary map
\[
\F(\C,\QF\qshift{1}) \to \Sig\F(\C^{\natural},\QF)
\]
arising from the metabolic sequence.
To show statement~\reffouritem{item:sig-component} we show momentarily that the composite
\[
(\C^\natural,\QF\qshift{1}) \to (\C,\QF)_\I^\flat \xrightarrow{\ev_{-1}} (\C^\natural,\QF\qshift{1})
\]
is given by the Poincar\'e functor $(\id,\eps(c)\cdot \id)$. Taking this for granted for the moment, the result follows immediately in case $\eps(c) =1$. If instead $\eps(c) = -1$, we recall from Corollary~\reftwo{corollary:inversion} that $\F(\id,-\id) = -\id_{\F(\C^\natural,\QF\qshift{1})} - \F(\mathrm{N}\Omega)$, where $\mathrm{N}\Omega$ is the norm, in the sense of Construction~\reftwo{construction:norm}, of the exact functor $\Omega$ on $\C^\natural$. Concretely, $\mathrm{N}\Omega$ is given by the composite
\[
(\C^\natural,\QF\qshift{1}) \xrightarrow{\Omega^\hyp} \Hyp(\C^\natural) \to (\C^\natural,\QF\qshift{1}).
\]
where the latter functor is the canonical functor using that $\Hyp(-)$ is left adjoint to the forgetful functor. This functor in turn factors as 
\[
\Hyp(\C^\natural) \xrightarrow{\mathrm{can}} \Met(\C^\natural,\QF\qshift{1}) \to (\C^\natural,\QF\qshift{1})
\]
whose second map is the Poincar\'e-Verdier projection of the metabolic sequence.
As a result, the composite 
\[
\F(\C^\natural,\QF\qshift{1}) \xrightarrow{\mathrm{N}\Omega} \F(\C^\natural,\QF\qshift{1}) \xrightarrow{\partial} \Sigma \F(\C^\natural,\QF)
\]
vanishes canonically, and the result follows.

Now, the remaining claim is implied by the statement that the
composite
\[
(\Spaf,\QF^\uni) \to (\Spaf,(\QF^\uni)\qshift{-1})_\I^\flat \xrightarrow{\ev_{-1}} (\Spaf,\QF^\uni)
\]
classifies $(\SS,\eps(c)\cdot \beta_u)$. To proceed we now have to make explicit the construction of $\bar{\sigma}^\vis_c(\I)$.
To that end, we note that the constant diagram $[\begin{tikzcd} \SS & \SS \ar[l,shift left,"\id"] \ar[l, shift right,"\id"'] \end{tikzcd}]$
maps to 
the constant diagram $\SS_{\B\ZZ}$ under the functor $\Spaf_\I \to \Spaf_{\B\ZZ}$, since $\I \to \B\ZZ$ is a localisation. %
Moreover, the induced map
\[
(\QF^{\uni})\qshift{-1}_\I(\SS \leftleftarrows \SS) =  \colim_\I (\QF^{\uni})\qshift{-1}(\SS) \to  \colim_{\B\ZZ}(\QF^{\uni})\qshift{-1}(\SS) = (\QF^{\uni}_{\B\ZZ})\qshift{-1}(\SS_X)
\]
is then an equivalence since any localisation is cofinal. %
We may consequently (essentially uniquely) lift the visible signature form $\bet_{c}$ on $\SS_{\B\ZZ}$ to a hermitian form $\bet^\I_c \in (\QF^{\uni})\qshift{-1}_\I([\SS \leftleftarrows \SS])$. Concretely, $\beta^\I_c$ is given by the composite
\begin{equation}
\label{equation:beta-I-c}%
\SS \xrightarrow{\bet_{\uni}} \QF^{\uni}(\SS) \xrightarrow{c\otimes \QF^{\uni}(\SS)} \SS\qshift{-1}[\B\ZZ] \otimes \QF^{\uni}(\SS) = \colim_{\B\ZZ} (\QF^{\uni})\qshift{-1}(\SS) \simeq \colim_{\I} (\QF^{\uni})\qshift{-1}(\SS),
\end{equation}
where $\bet_{\uni}$ is the universal Poincaré form on $\SS$. Given this, it is readily computed that the Poincar\'e functor $\ev_{-1}$ takes $\beta_c^\I$ to $\eps(c)\cdot \beta_u$.
It remains to note that the hermitian object $(\SS \leftleftarrows \SS,\bet^\I_c)$ in $(\Spaf,(\QF^{\uni})\qshift{-1})_\I$ is in fact Poincar\'e. For this, one can employ the embedding $\Spaf_\I \subseteq \Fun(\I\op,\Pro(\Spaf))$, which takes $\SS \leftleftarrows \SS$ to the constant diagram on $\SS$ and observe that also its dual is constant by the explicit formula \eqreffour{equation:explicit-dual-I}. But under the analogous embedding for $\Spaf_{\B\ZZ}$ the functor $p_* \colon \Spaf_\I \rightarrow \Spaf_{\B\ZZ}$ is induced by right Kan extension along $p\colon \I\op \to \B\ZZ\op$, which has restriction as a fully faithful adjoint. Since $\SS \leftleftarrows \SS$ and its dual both lie in the image of this restriction functor and the map connecting them is taken to the equivalence $\beta_c$ by $p_*$, it is itself an equivalence as claimed.
\end{proof}



We finish this subsection by discussing the difference between the geometric and algebraic decompositions of $\F(\C,\QF)_{\B\ZZ}$ into a direct sum of $\F(\C,\QF)$ and $\Sig\F(\C^{\natural},\QF)$, for a bordism invariant and Verdier-localising functor $\F$. The above shows that under the canonical equivalence $\Sigma \F(\C,\QF) \simeq \F(\C,\QF\qshift{1})$ the composite
\[
\begin{tikzcd}[column sep=large]
	\F(\C,\QF) \oplus \Sigma \F(\C^\natural,\QF) \ar[r,"\simeq"] & \F((\C,\QF)_{\B\ZZ}) \ar[r,"\simeq","{\reffour{corollary:algebraic-shaneson-splitting}}"'] & \F(\C,\QF) \oplus \Sigma \F(\C^\natural,\QF) 
\end{tikzcd}
\]
where the first and second equivalence are the ones from Corollaries~\reffour{corollary:geometric-shaneson-splitting} and~\reffour{corollary:algebraic-shaneson-splitting}, respectively,
is given by the following matrix:
\[
\begin{pmatrix}
	\id & 0 \\ \F(\tr^{\flat,q}_{\B\ZZ,c}) & \eps(c)\cdot \id
\end{pmatrix}
\]
Now, by construction, $\tr^{\flat,q}_{\B\ZZ,c}$ is given by multiplication by $q_*\vsign_c(\B\ZZ)$ and if the spherical fundamental class $c$ has Pontryagin-Thom class $q_*[c] = 0$, then $q_*(\vsign_c(\B\ZZ))$ is cobordant to $0$ in $(\Spaf,(\QF^\uni)\qshift{-1})$ by Remark~\reffour{remark:pontryagin}. Note that it is not evidently also cobordant to $0$ in the minimalisation of $\Spaf$, in fact, we will show later that it is \emph{not} null-cobordant therein. As a consequence of the above, we nevertheless obtain:

\begin{corollary} 
\label{cor:geometric=algebraic-decomposition}%
If the spherical fundamental class $c$ satisfies $q_*[c]=0$, the map 
\[
\F(\tr^{\flat,q}_{\B\ZZ,c})\colon \F(\C^\natural,\QF\qshift{1}) \to \F(\C,\QF)
\]
is null-homotopic after pre-composition with to $\F(\C, \QF\qshift{1})\to \F(\C^\natural,\QF\qshift{1})$ and after post-composition with $\F(\C,\QF) \to \F(\C^\natural,\QF)$. In particular, for idempotent complete Poincar\'e $\infty$-categories and for a spherical fundamental class $c$ with $(\eps(c),q_*[c]) = (1,0)$ in $\{\pm1\} \times \ZZ/2$, the algebraic and geometric decompositions of $\F(\C,\QF)_{\B\ZZ}$ coincide.
\end{corollary} 



Let us now investigate further the case where $(\eps(c),q_*[c]) = (1,0)$ but $\C$ is not necessarily idempotent complete. For this, let $(\C,\QF)\subseteq (\D,\QFD)$ be a Poincar\'e-Karoubi equivalence, so that $\QF = \QFD|_{\C}$ and $\C$ contains the minimal dense subcategory $\D^0\subseteq \D$ spanned by the objects with trivial K-class. Our goal is to give a more explicit description of the associated extended framed transfer
\[
-\otimes q_*\vsign_c(\B\ZZ)\colon \F(\D, \QFD\qshift 1) \to \F(\C, \QF)
\]
As a consequence of Corollary~\reffour{cor:geometric=algebraic-decomposition}, this map is canonically nullhomotopic if we restrict the source to $\F(\C, \QF\qshift 1)$ and if we map the target into $\F(\D, \QFD)$. This simple observation admits the following strengthening:



\begin{proposition}
\label{proposition:tau-is-equiv}%
In the above situation, the extended framed transfer factors canonically as 
\[
\F(\D,\QFD\qshift{1}) \to \cof[\F(\C,\QF\qshift{1}) \to \F(\D,\QFD\qshift{1})]
\xrightarrow\simeq  \fib[\F(\C,\QF) \to \F(\D,\QFD)] \to \F(\C,\QF).
\]
where the  outer maps are the canonical ones, and the middle one is an equivalence. 
\end{proposition}

\begin{proof}
Write $\Met(\D,\C,\QFD) \subseteq \Met(\D,\QFD)$ for the dense Poincaré subcategory spanned by those $\cob \to x$ such that $x \in \C$, and $\Ar(\D,\C,\QFD) \subseteq \Ar(\D,\QFD)$ for the dense Poincaré subcategory spanned by those $f\colon y \to z$ such that $\cof(f) \in \C$. We note that $\Met(\D,\C,\QFD) \simeq \Ar(\D,\C,\Phi\qshift{-1})$ via the 
equivalence of Lemma~\refone{lemma:arrow-is-met}. In addition, these Poincaré $\infty$-categories then fit into the solid commutative squares 
\[
\begin{tikzcd}
(\C,\QF\qshift{1}) \ar[r]\ar[d,"\id_{(-)}"'] 
  & (\D,\QFD\qshift{1}) \ar[d,"\id_{(-)}"]  
  && \Met(\D,\C,\QFD) \ar[d,"\met"']\ar[r] 
  & \Met(\D,\QFD) \ar[d,"\met"] 
  \\
  \Ar(\C,\QF\qshift{1}) \ar[r] 
  & \Ar(\D,\C,\QFD\qshift{1})\ar[urr,dotted,"{-\otimes z}"] && (\C,\QF) \ar[r] & (\D,\QFD)
\end{tikzcd}
\]
in which the solid horizontal arrows are dense Poincaré-Karoubi inclusions. We now observe that the right solid square in the diagram is cartesian by construction and its right vertical arrow is a split Poincaré-Verdier projection, so that the whole square is a split Poincaré-Verdier square by Proposition~\reftwo{corollary:PV-proj-pullback}.
On the left hand side we have the dual phenomenon: the vertical arrows are both split Poincaré-Verdier inclusions (with left and right adjoint sending an arrow to its target and source, respecitvely), and the square is vertically left and right adjointable by direct inspection. The induced functor on vertical cofibres is an equivalence (where both sides can be identified with $(\C,\QF)$ via $\cof$), and so by Proposition~\reftwo{proposition:poincare-verdier-square-char} we conclude that the left solid square is cocartesian. 
The functor $\F$ therefore takes both solid squares to cartesian squares by additivity, and since $\F$ vanishes on $\Ar(\C,\QF)$ and $\Met(\C,\QF)$ by bordism invariance, we obtain identifications
\[
\F(\Ar(\D,\C,\QFD\qshift{1})) = \cof[\F(\C,\QF\qshift{1}) \to \F(\D,\QFD\qshift{1})] \quad\text{and}\quad \F(\Met(\D,\C,\QFD)) = \fib[\F(\C,\QF) \to \F(\D,\QFD)] .
\]


We next construct the diagonal Poincar\'e functor marked by a dotted arrow in the diagram: For this, we identify its target with $\Ar(\D, \C, \QF\qshift{-1})$ as above and will use that $\Ar\colon \Catp\to \Catp$ is lax symmetric monoidal for the tensor product of Poincar\'e categories, as a composite of the strong symmetric monoidal inclusion $\Catp\to \Cath$ and the lax symmetric monoidal right adjoint $\Pairings\colon \Cath\to \Catp$. On underlying categories, this lax symmetric monoidal structure is pointwise, i.e.\ given by 
\[
\Ar(\C)\otimes \Ar(\D)\to \Ar(\C\otimes \D), \quad (c\to c')\otimes (d\to d') \mapsto (c\otimes d\to c'\otimes d').
\]
By passing to dense subcategories, this symmetric monoidal refinement of $\Ar(-)$ induces a Poincar\'e functor 
\[
\Ar(\D, \C, \QFD\qshift n)\otimes \Ar(\Spaf, \Spa^0, (\QF^{\uni})\qshift m)\to \Ar(\D, \C, \QFD\qshift{n+m}),
\]
where we write $\Spa^0$ for the minimalisation $(\Spaf)^{\min}$ of $\Spaf$ consisting of all objects with trivial $\K$-class. 


Let us now fix a null-homotopy of Pontryagin-Thom map %
$q_*[c]\colon \SS \to \SS\qshift{-1}$,
which by Remark~\reffour{remark:pontryagin}
determines %
a Lagrangian $(\SS \to \SS^{\B\ZZ},\bet_c,\eta)$, and let us write 
\[
z \in \Poinc(\Met(\Spaf, \Spa^0, (\QF^{\uni})\qshift{-1}))\simeq \Poinc(\Ar(\Spaf, \Spa^0, (\QF^{\uni})\qshift{-2}))
\]
for this corresponding Poincaré object.
In particular, the underlying object in $\Met(\Spaf)$ is given by the constant inclusion $\SS\to \SS^{\B\ZZ}$, and the corresponding object in $\Ar(\Spaf)$ is obtained by passing to fibres, giving a map $\SS\qshift{-2}\to \SS$. 
The diagonal Poincar\'e functor $\Ar(\D, \C, \QF\qshift{1}) \to \Ar(\D, \C, \QF\qshift{-1})$ in the above diagram is defined by tensoring with this Poincar\'e object $z = (\SS\qshift{-2} \to \SS,\gam)$. The composite functor
\[
(\D,\QFD\qshift{1})  \to \Ar(\D,\C,\QFD\qshift{1}) \xrightarrow{(-)\otimes z} \Ar(\D,\C,\QFD\qshift{1}) \xrightarrow{\cof} (\C,\QF)
\]
is then given by tensoring with $q_*\vsign_c(\B\ZZ)$.


We can therefore finish the proof by showing that the diagonal functor is a bordism equivalence. 
For this, it will suffice to show that the Poincaré object $z$
is graded bordism invertible, that is, it admits a tensor inverse in $\Poinc(\Ar(\Spaf,\Spa^0,(\QF^{\uni})\qshift{2}))$ up to cobordism. Let $(\SS^2,0) \in \spsforms(\Spaf,(\QF^{\uni})\qshift{2})$ be the 2-shifted hermitian object whose underlying object is the $2$-shifted sphere spectrum and whose hermitian form is $0$. Under proposition Proposition~\refone{proposition:thom} this corresponds to a hyperbolic Poincaré object of the form $([\SS^2 \xrightarrow{0} \SS],\hyp) \in \Poinc(\Ar(\Spaf,(\QF^{\uni})\qshift{2}))$, and since $\SS^2$ and $\SS$ have the same $\K$-class this Poincaré object actually lies in $\Ar(\Spaf,\Spa^0,(\QF^{\uni})\qshift{2})$. Its tensor product with $z=(\SS\qshift{-2}\to \SS, \gam)$ is
\[
([\SS^2 \xrightarrow{0} \SS],\hyp) \otimes ([\SS^{-2} \to \SS],\gam) = ([\SS \xrightarrow{0} \SS],\hyp),
\]
the hyperbolic Poincaré object corresponding to the trivial (unshifted) form on $\SS$. We now claim
that $([\SS \xrightarrow{0} \SS],\hyp)$ is cobordant to the unit Poincaré object $([\SS \xrightarrow{\id} \SS],u)$ of $\Ar(\Spaf,\Spa^0,\QF^{\uni})$. Indeed, this follows from the cofinality theorem \reffour{theorem:strong-cofinality}, whose $\L$-theory version gives 
\[
\L_0(\Ar(\Spaf, \Spa^0, \QF^{\uni}) \simeq \pi_0 (\L(\Spaf, \QF^{\uni})/\L(\Spa^0, \QF^{\uni})) \simeq \pi_0 (\K_0(\Spaf)/\K_0(\Spa^0))^\tC \simeq \ZZ/2
\]
where the composite map explicitly sends $c\to d$ to the $\K_0$-class of $d$. 
\end{proof}




 
Let us finally examine the special case where $\F$ is the $\L$-spectrum functor. 
By Theorem~\reffour{theorem:strong-cofinality}
the equivalence of Proposition~\reffour{proposition:tau-is-equiv} then takes the form
\[
\Sig[\K_0(\D)/\K_0(\C)]^{\tC} \to \Om[\K_0(\D)/\K_0(\C)]^{\tC} .
\]
which on $\pi_{-n}$ is the canonical isomorphism in Tate cohomology
\[
\hat{\mathrm{H}}^{n+1}(\Ct,\K_0(\D)/\K_0(\C)) \xrightarrow{\cong} \hat{\mathrm{H}}^{n-1}(\Ct,\K_0(\D)/\K_0(\C)) .
\]
Indeed, by naturality and the classification of dense Poincaré subcategories, it is enough to note that if $\K_0(\C,\QF)$ is cyclic then all its Tate cohomologies are
either $0$ or $\ZZ/2$, 
and in any case admits no non-identity automorphisms (compare Remark~\reffour{remark:Rothenberg_map}).
We may hence view Proposition~\reffour{proposition:tau-is-equiv} as a generalisation of Ranicki's \cite{Ranicki-splitting}*{Proposition 6.2}, where this statement is proven specifically for quadratic L-theory of group rings. 



\begin{example}
\label{example:concrete-difference}%
Suppose that $R$ is a number ring in which 2 is invertible and such that $\Pic(R)/2$ is non-trivial. For example, we may take $R = \ZZ[\sqrt{-pq},\tfrac{1}{2}]$ where $p$ and $q$ are two distinct positive odd primes; the 2-rank of $\Pic(\ZZ[\sqrt{-pq}])$ is then $2$ by Gauss' genus formula (see, e.g., \cite{frohlich-taylor}*{\S V.I, Theorem 39}), and so the 2-rank of $\Pic(R)$ is at least 1. %
As explained in Example~\reffour{example:free-dedekind}, in this case the map
\[
\L_{2k+1}(\Dperf(R),\QF^{\sym}_R) \to \pi_{2k+1}[\K_0(\Dperf(R))/\K_0(\Dfree(R))]^{\tC} = \Pic(R)/2
\]
is an isomorphism when $k$ is odd and zero when $k$ is even (since the assumption that $2$ is invertible in $R$ implies that $R$ has no dyadic primes).
It then follows that the composite
\[
\begin{tikzcd}
\L_{-1}(\Dperf(R),\QF^{\sym}_R) \ar[r, ->>] & \Pic(R)/2 \ar[rrr,"\tau_{(\Dperf(R),\Dfree(R),\QF^{\sym})}","\simeq"'] &&& \Pic(R)/2 \ar[r,"\partial", hook] & \L_0(\Dfree(R), \QF^{\sym}_R)
\end{tikzcd}
\]
is non-zero as soon as $\Pic(R)/2$ is non-zero. In particular, in this case the geometric and algebraic decompositions of $\L((\Dfree(R),\QF^{\sym}_R)_{\B\ZZ})$ into a direct sum $\L(\Dfree(R),\QF^{\sym})$ and $\Sig\L(\Dperf(R),\QF^{\sym})$ do not coincide.
\end{example}
